\documentclass[11pt]{article}

\usepackage[T1]{fontenc}
\usepackage[utf8]{inputenc}
\usepackage{lmodern}
\usepackage[margin=1in]{geometry}
\usepackage{amsmath,amssymb,amsthm,mathtools}
\usepackage{enumitem}
\usepackage{booktabs}
\usepackage{array}
\usepackage{longtable}
\usepackage{xcolor}
\usepackage{microtype}
\usepackage{etoolbox}
\usepackage{url}
\usepackage[colorlinks=true,linkcolor=blue!50!black,citecolor=green!35!black,urlcolor=blue!60!black]{hyperref}
\usepackage[capitalize,nameinlink]{cleveref}

% ---------------------------------------------------------------------------
% Theorem environments
% ---------------------------------------------------------------------------
\theoremstyle{plain}
\newtheorem{theorem}{Theorem}[section]
\newtheorem{lemma}[theorem]{Lemma}
\newtheorem{corollary}[theorem]{Corollary}
\newtheorem{proposition}[theorem]{Proposition}
\newtheorem{conjecture}[theorem]{Conjecture}
\newtheorem{question}[theorem]{Question}
\newtheorem{problem}[theorem]{Problem}
\theoremstyle{definition}
\newtheorem{definition}[theorem]{Definition}
\newtheorem{remark}[theorem]{Remark}
\newtheorem{example}[theorem]{Example}

% The Lean proof is the proof; the text gives a faithful outline of it.
\newenvironment{outline}[1][Outline of the formal proof]{\begin{proof}[#1]}{\end{proof}}
% Results that are not formalized get complete proofs, labelled as such.
\newenvironment{informalproof}{\begin{proof}[Proof (not formalized)]}{\end{proof}}

% ---------------------------------------------------------------------------
% Lean references
%   \lean{name}            a declaration, relative to the namespace NativeRelations
%   \formal{file}{names}   the formal counterpart of a numbered result: a file relative to
%                          lean/NativeRelations/ (or, when every declaration belongs to an earlier
%                          library, a file of that library) and a comma-separated list of declarations;
%                          names from earlier libraries carry their namespace prefix
%   \partformal{file}{names}{gap}
%                          a numbered result whose formal counterpart covers all but the
%                          part described in {gap}; the gap is proved or cited in the text
%   \notformalized         a numbered result that has no formal counterpart
%   \cited{where}          a result of an earlier paper used without proof: where it is proved,
%                          and whether it is Lean-verified there
% ---------------------------------------------------------------------------
% Lean names are typeset like paths, with line breaks allowed after "_", "." and "/".
\DeclareUrlCommand\leanname{\def\UrlFont{\ttfamily\upshape}}
\renewcommand{\UrlBreaks}{\do\_\do\.\do\/\do\-}
\newcommand{\lean}[1]{\leanname{#1}}
\newcommand{\leanfile}[1]{\leanname{#1}}
\newcommand{\leanlist}[1]{\begingroup\def\leansep{}%
  \renewcommand*{\do}[1]{\leansep\leanname{##1}\def\leansep{, }}\docsvlist{#1}\endgroup}
\newcommand{\formal}[2]{\par\noindent{\small\upshape\raggedright\textsf{Lean:} \leanlist{#2}\unskip
  \nobreak\hfil\penalty50\hskip1em\hbox{}\nobreak\hfill[\leanfile{#1}]\parfillskip=0pt\par}%
  \smallskip\noindent\ignorespaces}
\newcommand{\partformal}[3]{\par\noindent{\small\upshape\raggedright\textsf{Lean (partial):} \leanlist{#2}\unskip
  \nobreak\hfil\penalty50\hskip1em\hbox{}\nobreak\hfill[\leanfile{#1}]\parfillskip=0pt\par}%
  \par\noindent{\small\upshape\textsf{Not formalized:} #3}\par\smallskip\noindent\ignorespaces}
\newcommand{\notformalized}{\par\noindent{\small\upshape\textsf{Not formalized;} a complete proof is given
  below.}\par\smallskip\noindent\ignorespaces}
\newcommand{\cited}[1]{\par\noindent{\small\upshape\textsf{Cited:} #1}\par\smallskip\noindent\ignorespaces}

% ---------------------------------------------------------------------------
% Notation
% ---------------------------------------------------------------------------
\newcommand{\N}{\mathbb{N}}
\newcommand{\Z}{\mathbb{Z}}
\newcommand{\Q}{\mathbb{Q}}
\newcommand{\R}{\mathbb{R}}
\newcommand{\C}{\mathbb{C}}
\newcommand{\F}{\mathbb{F}}
\newcommand{\T}{\mathsf{T}}
\newcommand{\clog}[1]{\lceil \log_2 #1 \rceil}
\newcommand{\abs}[1]{\lvert #1 \rvert}
\newcommand{\card}[1]{\lvert #1 \rvert}
\newcommand{\eps}{\varepsilon}
\newcommand{\prev}{\mathrm{prev}}
\newcommand{\cost}{\mathrm{cost}}
\newcommand{\gn}{\mathrm{gn}}
\newcommand{\Sh}{\mathrm{Sh}}
\newcommand{\cl}{\mathrm{cl}}
\newcommand{\Aut}{\mathrm{Aut}}
\newcommand{\minrk}{\mathrm{minrk}}
\newcommand{\MAIS}{\mathrm{MAIS}}
\newcommand{\Crux}{\mathrm{Crux}}
\newcommand{\TCP}{\mathrm{TCP}}
\DeclareMathOperator{\rank}{rank}
\DeclareMathOperator{\dist}{dist}
\DeclareMathOperator{\poly}{poly}

% ---------------------------------------------------------------------------
% Notation of this paper
% ---------------------------------------------------------------------------
\newcommand{\wval}{w_{\mathrm{val}}}
\newcommand{\Lin}{\mathsf{In}}
\newcommand{\Out}{u_0}
\newcommand{\CL}{\mathrm{CL}}
\newcommand{\Hop}{\mathrm{Hop}}
\newcommand{\tpos}{\mathrm{tpos}}
\newcommand{\cE}{c_{\mathrm{E}}}
\newcommand{\thetaM}{\theta_{\mathrm{mac}}}
\newcommand{\Mmax}{M_{\max}}
\newcommand{\Ns}{N_s}
\newcommand{\Qstar}{Q^{*}}
\newcommand{\Range}[1]{[0,#1)}


\title{Oblivious Linear-Time Tape Transposers\\ Need Macroscopic Crossings}
\author{Hans Brende \and Claude Opus 5.5 (Anthropic) \and Gemini 4 Argon (Google DeepMind)}
\date{November 2026}

\begin{document}

\maketitle

\begin{abstract}
Can a multitape Turing machine transpose a $d\times d$ matrix of $\lceil\log_2 d\rceil$-bit entries in time linear in its
length~$L$? We show that an oblivious machine that does so must route a linear amount of information through
\emph{macroscopic crossings}: pairs of revisits of one tape, each taking the head a fixed fraction of~$L$ away from its cell,
that interleave as a queue's would. Precisely, for every set of at least $\alpha L$ output bits, every set of steps that separates
their output steps from the earlier ends of such crossings in the time--space dependency graph has at least $c$ times as many
elements as the output set. Here $c$ and the fraction depend only on the number of tapes, the alphabet, the number of states, the
time constant and~$\alpha$. The expansion can be found on a single coarsely monotone leg of a single tape, in a coarse skeleton of the
run drawn from a finite family independent of~$d$, and at least two work heads must make revisits of linear excursion. Behind these
statements is a dichotomy for arbitrary linear-time traces: either such crossings are reached with linear expansion, or a separator
exists that is cheaper than the acyclic set it leaves in the cone relation, which no transposer can have. For adaptive machines the
conclusions hold on every trace followed by an exponentially small but fixed fraction of the inputs, and on the heavy traces of a bit
simulator, with the same tapes and alphabet, of every linear-time transposer; carrying them back to the machine's own traces is open.
We also prove that every cut between a set of input rows and the outputs they feed into a set of output rows carries the corresponding
information whatever the coding, and we exclude classes of machines with boundedly many macroscopic phases, oblivious and adaptive.
Explicit linear-time schedules, realized by machines that compute other functions, meet these conditions in every test we could run
(numerically, at $d\le128$), so the conditions are not known to decide the question. We show that one set of steps separates the
inputs from the outputs for every pair of row sets at once, so counting such cuts cannot by itself go beyond linear time, and we state
the multi-scale question that remains. The main theorems are proved in Lean~4; one proposition and one corollary are proved in the
text, and the exclusions of individual schedule families are given as proof sketches. The constants of the main theorems are explicit
and astronomically small; those of the two-heads theorem are polynomial, and those of one class exclusion come from a compactness
argument.

\end{abstract}

\begin{quote}\small
\textbf{Contributions.} Hans Brende chose the problem, fixed the verification standard, and directed the research.
Gemini 4 Argon, a large language model developed by Google DeepMind, developed most of the mathematics and wrote the
Lean formalization. Claude Opus 5.5, a large language model developed by Anthropic, reviewed the mathematics and the Lean
throughout, refereed the main proofs, proposed several of the statements proved here (among them the macroscopic form of
the crossing theorem and the trace-only dichotomy), and wrote this paper under Hans Brende's direction.
\end{quote}

\tableofcontents
\clearpage

\section{Introduction}\label{sec:intro}

\subsection{The problem}

Can a deterministic multitape Turing machine transpose a $d\times d$ matrix with $b(d)=\max(1,\clog d)$-bit entries, written
row by row, in time linear in the input length $L=d^2b(d)$? This is the statement~(T) of the first paper of this
series~\cite{FirstPaper2026}. By a reduction of Harvey and van der Hoeven~\cite[Cor.~6.9]{HvdH2025}, a linear-time
multiplication machine would give a linear-time transposer, so a proof of $\neg$(T) would give the first superlinear lower
bound for off-line integer multiplication on multitape Turing machines. The question is open, and this paper does not decide
it.

A machine is \emph{oblivious} if its head motion and the steps at which it appends output bits depend only on~$d$. Every machine can
be made oblivious at a logarithmic cost in time~\cite{PippengerFischer1979}, so oblivious machines are the natural first target. The
run of an oblivious machine is a fixed time--space network, and a linear-time oblivious transposer is a network code for the
transposition demands on that network. Adler, Harvey, Jain, Kleinberg and Lehman showed that if coding could not beat routing on
undirected networks, oblivious two-tape transposition would need time $\Omega(L\log d)$~\cite[\S6]{AHJKL2006}; Braverman and He, and
concurrently Zhang, Li and Li, have refuted that hypothesis~\cite{BravermanHe2026,ZhangLiLi2026}, so for oblivious machines the question
is genuinely one about codes. The series has reduced it to codes in two steps. The fourth paper reduced $\neg$(T) for oblivious machines
to depth reduction and a coding statement, the \emph{crux}~\cite[Theorem~D]{FourthPaper2026}. Lower bounds for several restricted
classes of labels suggested that depth reduction might fail~\cite{FourthPaper2026,EighthPaper2026,DirectedReciprocity2026,MergeDepth2026},
but the fourteenth paper proved it for every number of heads, so that the crux alone would exclude oblivious linear-time
transposers~\cite[Corollary~O]{BoundedStack2026}. The crux asks whether the \emph{cone relations} of a linear-time schedule, the
relations between each output bit and the input bits it can see without crossing a small revealed set, admit codes of rate close to~$1$.
Machine-free relations of this kind carry storage codes above rate one half~\cite{StorageCodes2026}, and binary codes of rate approaching
one exist on suitable separated relations~\cite{HighRate2026}; whether the relations of an actual linear-time schedule can carry them is
the question.

This paper attacks the crux from the side of the schedule. Instead of a coding bound for all relations, we ask which schedules make the
crux \emph{provable}. We show that it holds, with an explicit deficit, for every schedule in which a small enough set of steps separates a
large set of outputs from the \emph{macroscopic crossings} of the tapes. An oblivious linear-time transposer must therefore route a linear
amount of information through macroscopic crossings, and we describe where it must do so.

\subsection{Results}

The run of a machine defines a \emph{dependency DAG} on its steps $0,\dots,T-1$: every step depends on the step before it, which
carries the state, and on the previous visit of every work head to the cell it is scanning, which carries that cell's
contents. A \emph{revisit} of a work tape is a pair $u<t$ of consecutive visits of its head to one cell between which the head leaves the
cell; its \emph{excursion} is the farthest the head travels from the cell before returning. Two revisits $u<t$ and $u'<t'$ of the same tape
\emph{cross} if $u<u'<t<t'$. A unit-speed head makes crossing revisits only when it reuses a stretch of tape like a queue, across two
turns: write a stretch moving right, return, and read it again moving right. Reuse like a stack, across one turn, makes only nested
revisits. Given $\theta>0$, a step is a \emph{macroscopic crossing step} if it is the earlier end of a revisit of excursion more than
$\theta L$ that crosses another revisit of excursion more than $\theta L$ on the same tape. A trace of a machine is \emph{$\eta$-heavy} if at
least a $2^{-\eta L}$ fraction of the inputs follow it; an oblivious machine has one trace, and it is $0$-heavy. \Cref{sec:prelim} states
these definitions exactly.

\begin{itemize}[leftmargin=1.5em]
\item \textbf{Theorem~X, macroscopic crossings} (\cref{thm:X}). Let $M$ be an oblivious machine with $k$ work tapes that
transposes $d\times d$ matrices within $C(L+1)$ steps, and let $\alpha>0$. There are explicit $\theta>0$, $c>0$ and $d_0$,
depending only on $k$, the alphabet and state counts of~$M$, $C$ and~$\alpha$, such that for every $d\ge d_0$ the following
holds. For every set $V_0$ of at least $\alpha L$ output positions, every set of steps that meets every dependency path from
the output steps of $V_0$ to the macroscopic crossing steps has at least $c\abs{V_0}$ elements. By Menger's theorem, there are then
$c\abs{V_0}$ vertex-disjoint dependency paths from these outputs into macroscopic crossings.
\item \textbf{Theorem~D, a dichotomy for all traces} (\cref{thm:D}). Theorem~X is the transposer case of a statement about
arbitrary linear-time traces, with no hypothesis on what the machine computes: for every output set of size at least $\alpha L$,
either the macroscopic crossings are reached with linear expansion, or a cheap separator leaves a linear-size acyclic set in the cone
relation. A separator cheaper than the acyclic set it leaves, at a radius at least the genie depth, is impossible for a transposer
(\cref{lem:acyclic-cost}).
\item \textbf{Theorem~S, skeletons} (\cref{thm:S}). At scale $\theta L/4$ the walk of each work tape of a linear-time run is a
word of at most $\Mmax=\lfloor16C/\theta\rfloor$ steps $\pm1$, so the coarse shapes of all runs come from an explicit finite
family of $(2^{\Mmax+1}-1)^k$ tuples, independent of~$d$. For an oblivious linear-time transposer the shape is
\emph{crossing-capable} (some tape turns at least twice at this scale with overlapping legs), and for every large output set the
expansion of Theorem~X is carried by the crossing steps of a \emph{single coarsely monotone leg of a single tape}. Along infinitely many
sizes, one fixed shape, tape and leg serve for the set of all input positions.
\item \textbf{Theorem~H, two long-range heads} (\cref{thm:H}). No machine that transposes in linear time with an $\eta$-heavy trace, for a
small fixed $\eta$, has on that trace one unrestricted work head while every other work head keeps its revisits within a small fixed
fraction of~$L$, whatever the input head does. Every such transposer, in particular every oblivious one, therefore has at least two work
heads that return to cells after moving a fixed fraction of $L$ away; and for each work head, some set containing a fixed fraction of its
outputs reaches long revisits off that head with linear expansion.
\item \textbf{Theorem~E, adaptive machines} (\cref{thm:E}). The conclusion of Theorem~X holds for every $\eta$-heavy trace of an arbitrary
linear-time transposer, for an explicit $\eta>0$. It holds on all heavy traces of the bit simulator of the fifteenth
paper~\cite[Theorem~S]{HighRate2026} applied to any linear-time transposer; that simulator has the same tapes and work alphabet, runs in
linear time, and its heavy traces carry all but a $2^{-\eta L}$ fraction of the inputs. Transferring the conclusion back to the original
machine's own dependency DAG is open; every statement in this paper about the dependency DAG of a general machine assumes a heavy trace or
names the simulator.
\item \textbf{Lemma~R, rows to columns} (\cref{lem:R}). For every machine that transposes at size $d$, every $\eta$-heavy trace, and every two
sets $A_0,A_1$ of rows, every set of steps that separates the input reads of the rows in $A_0$ from the output steps of the entries that go
from rows $A_0$ to output rows $A_1$ has at least $(\abs{A_0}\abs{A_1}b-\eta L)/(\wval+1)$ elements, where $\wval$ is the number of bits of
one step's value.
\item \textbf{Class exclusions} (\cref{sec:classes}). Machines whose heads make boundedly many macroscopic phases, oblivious or adaptive with
sublinear jitter, are excluded by an explicit native deficit (Theorems~B and~J); a fixed native deficit suffices in place of the half-rate
hypothesis of the fourteenth paper's low-adaptivity corollary (Theorem~G); and machines with no free work head, bounded nesting load and
boundedly many scales are excluded by a scale-gap pigeonhole (Theorem~Q). \Cref{app:catalogue} lists explicit schedule families excluded
by these theorems and by certificates, with proof sketches.
\end{itemize}

\paragraph{What remains.} Explicit oblivious linear-time schedules, realized in our experiments by machines that compute other functions,
meet the necessary conditions above in every test we could run (numerically, at $d\le128$; \cref{sec:construction}). So these conditions
are not known to decide the question. Counting cuts alone cannot go further: one set of steps, the storage live across the last input read
together with the earlier output steps, separates the inputs from the outputs for every pair of row sets at once (\cref{prop:onecut}).
Outside restricted schedule classes (\cref{cor:scale-separated}), the only coding-proof logarithmic factor known on tapes comes from counting
passes (Theorem~P of the first paper), and \cref{sec:open} states the multi-scale question it leaves open.

\paragraph{Constants.} The constants of Theorems~X, D and~S are explicit, positive and astronomically small. For two work tapes, time
constant $C=1$ and $\alpha=1$, the expansion constant of Theorem~X is $\rho_0/(4\wval)$ with $\rho_0\approx10^{-105}$, and the threshold $\theta$
is smaller still; \cref{sec:crossings} gives the formulas. The constants of Theorem~H are polynomial in the time constant and the value width,
and those of Theorem~Q come from a compactness argument and are not explicit. The small constants come from the proofs: on the schedules we
measured, the expansion ratio was exactly~$1$ at every tested threshold up to $L/2$ (numerically, at $d\le128$; \cref{rem:numerics-crossings}).

\subsection{How to read the paper and its formal proofs}

Every numbered result states its status. \textsf{Lean:} lines name the declarations of the accompanying Lean~4 library that
prove it, in the file named at the right; the Lean statement is authoritative and the text gives a faithful outline of the formal
proof. A result marked \textsf{Not formalized} has a complete written proof in the text. \textsf{Cited:} lines name results of
earlier papers of the series used without proof, and say whether they are Lean-verified there. Numerical observations appear
only in remarks or parentheses labelled as numerical and are never used in proofs. The formal library depends only on the axioms
\lean{propext}, \lean{Classical.choice} and \lean{Quot.sound} (\cref{app:verification}).

The formalization is in the frozen machine model of the first paper~\cite{FirstPaper2026}, through the traces, dependency DAGs
and index codes of the eleventh paper~\cite{MergeDepth2026}, the genies of the fourteenth~\cite{BoundedStack2026} and the
simulator of the fifteenth~\cite{HighRate2026}. Not formalized: \cref{prop:onecut} and \cref{cor:scale-separated}, which
have complete proofs in the text; the remarks marked as proof sketches; and the exclusions of the individual schedule families of
\cref{app:catalogue}, which are not numbered results and are given as proof sketches over the formal class theorems and certificate
lemmas.

\subsection{Related work}

\emph{Lower bounds on tapes.} Superlinear lower bounds for explicit problems are known for one work tape: Hennie's crossing
sequences~\cite{Hennie1965}, Maass's combinatorial arguments~\cite{Maass1985}, and for transposition itself the bounds of
Dietzfelbinger, Maass and Schnitger~\cite{DMS1991} and Dietzfelbinger and Maass~\cite{DietzfelbingerMaass1993} in the classical one-tape
models. With two or more work tapes the only superlinear lower bounds known are those that follow from the separation of deterministic
and nondeterministic linear time~\cite{PPST1983}, for problems that are hard for nondeterministic linear time~\cite{Grandjean1988}; none is
known for off-line multiplication, sorting, the discrete Fourier transform or transposition~\cite[\S1.4]{HvdH2025}. The simulation of time
by space~\cite{HPV1977} also rests on the structure of multitape computation graphs. For transposition, moving the entries as indivisible
items along a fixed pattern needs time $\Omega(L\log d)$~\cite{Stoss1973}, and similar bounds hold for indivisible records in external
memory~\cite{AggarwalVitter1988}; Kirchherr proved an $\Omega(\log d)$ bound on reversals for conservative machines~\cite{Kirchherr1988}, and
the first paper of this series an $\Omega(\log d)$ bound on monotone passes for arbitrary computation, with coding
allowed~\cite[Theorems~P and~P$'$]{FirstPaper2026}, extended to adaptive machines without a runtime assumption
in~\cite{PassesMemo2026}. The stream lower bounds of Grohe, Hernich and Schweikardt~\cite{GroheSchweikardt2005,GHS2009} bound reversals
for sorting-type problems with small internal memory. A linear-time machine, however, may reverse its heads $\Theta(L)$ times at small
scales, so pass and reversal bounds apply to it only after its small-scale motion is absorbed into the finite control.

\emph{Macroscopic motion.} The first paper's Theorem~M*~\cite[Theorems~7.8 and~7.9]{FirstPaper2026} passes to the limit of rescaled head
trajectories and shows that the heads of any linear-time transposer, oblivious or adaptive, satisfy an integral inequality forcing
returns to macroscopically earlier positions. The fourth paper deduced, by a written proof, that in every macroscopic limit of a linear-time
transposer at least two work-head limits are non-monotone~\cite[Corollaries~3.18 and~3.24]{FourthPaper2026}. Theorem~H(b) is a
quantitative form of this at a single size, with explicit constants and on heavy traces. Theorem~X asks for more than non-monotone motion: a
run whose heads reuse their tapes only like stacks has no macroscopic crossing steps, so Theorem~X excludes it whatever its macroscopic limit.

\emph{Network and index coding.} The time--space view of oblivious machines and the conditional lower bound are due to Adler et
al.~\cite{AHJKL2006}; the multiple-unicast conjecture was refuted by Braverman and He~\cite{BravermanHe2026} and, concurrently, by Zhang, Li
and Li~\cite{ZhangLiLi2026}. Our exclusions use the acyclic-set lower bound of index coding~\cite{BBJK2011}, in the machine form of the
eleventh paper~\cite{MergeDepth2026}.

\emph{Upper bounds.} Recursive block transposition takes time $O(L\log d)$. An unrefereed preprint claims time
$O(L(\log d)^{1-\kappa})$ for a small $\kappa>0$ through an address-interchange primitive that mixes data by exclusive
or~\cite{OpenAI2026mult}; if it is correct, coding already beats routing on tapes by a factor $(\log d)^{\kappa}$, and any
proof of $\neg$(T) must allow such gains.

\subsection{Organization}

\Cref{sec:prelim} sets out the model, the dependency DAG, revisits and crossings, cuts and cone relations. \Cref{sec:certificates}
gives the certificate method: an acyclic set of linear size in a cone relation, possibly after a cheap cut, excludes linear-time
transposition. \Cref{sec:heads} proves Theorem~H, \cref{sec:crossings} Theorems~D and~X, and \cref{sec:skeletons} Theorem~S.
\Cref{sec:demand} proves Lemma~R and the colouring criterion built on it. \Cref{sec:adaptive} treats adaptive machines,
\cref{sec:classes} the class exclusions. \Cref{sec:construction} turns the results around into a specification for a linear-time
transposer and describes the schedules that meet it, and \cref{sec:open} lists the open problems that would decide the question.
\Cref{app:catalogue} catalogues the excluded schedule families, \cref{app:leanindex} indexes the Lean declarations, and
\cref{app:verification} explains how to build and check the formalization.

\section{The model, traces and the dependency DAG}\label{sec:prelim}

\subsection{Machines and transposition}

We work in the frozen machine model of the first paper~\cite{FirstPaper2026}. A machine $M$ has a read-only two-way input
tape, $k=k(M)$ work tapes over the alphabet $\{0,\dots,g\}$ with blank~$0$, states $\{0,\dots,q\}$ with start state~$0$, and a
one-way write-only output tape that receives at most one bit per step. Every step writes a symbol on every work tape (possibly
the scanned one) and moves every head by at most one cell. The \emph{value width}
\[
  \wval(M)=\clog{(q+1)}+k\clog{(g+1)}+2
\]
bounds the number of bits needed to record what one step does: the new state, the $k$ written symbols and the appended output
bit. Always $\wval\ge2$.

For a side length $d\ge1$ put $b=b(d)=\max(1,\clog d)$, $W=db$ and $L=d^2b$. A $d\times d$ matrix with $b$-bit entries is
written row by row as a word of $L$ bits: the bit $\gamma$ of entry $(i,j)$ sits at position $(di+j)b+\gamma$. The transpose
has the same encoding with $i$ and $j$ exchanged, so output position $\tpos(q)$ carries input bit~$q$, where
\[
  \tpos\bigl((di+j)b+\gamma\bigr)=(dj+i)b+\gamma ,
\]
an involution of $\Range L$. For $q=(di+j)b+\gamma$ we call $i$ the \emph{input row} of $q$ and $j$ its \emph{output row}: the
copy of $q$ is written in row~$j$ of the output. $M$ \emph{transposes at size $d$ within time $c$}, written
$\mathrm{TransposesAt}(M,d,c)$, if on every matrix input it halts within $\lfloor c(L+1)\rfloor$ steps with the encoding of the
transpose on its output tape. Statement~(T) asks for a machine and a constant $c$ with this property at every~$d$.

$M$ is \emph{oblivious at length $L$ with halting time $T$} if every input of length $L$ halts at exactly step~$T$ and any two
such inputs have the same input-head position, the same work-head positions and the same append decision at every step. We
write $\mathrm{ObliviousLinTime}(M,d,C,T)$ if $M$ is oblivious at length $L$ with halting time $T\le C(L+1)$ and transposes at
size~$d$. This is a condition at a single size~$d$.

\subsection{Traces}

The \emph{trace} of $M$ on an input $w$ records, step by step, the move of the input head, the moves of the work heads and
whether a bit is appended~\cite{MergeDepth2026}. It determines the positions $y_j(t)$ of work head $j$ after $t$ steps, with
$y_j(0)=0$ and $\abs{y_j(t+1)-y_j(t)}\le1$, the input-head position, and the \emph{append steps} $u_0<u_1<\dots$ at which output
bits are written. For an oblivious machine all inputs of length $L$ share one trace, the trace on the all-zero input.

For a trace $\tau$ the \emph{fibre} of $\tau$ is the set of inputs $X\in\{0,1\}^L$ whose trace is~$\tau$, and the \emph{trace
probability} $\Pr(\tau)$ is the size of the fibre divided by $2^L$. A trace is \emph{$\eta$-heavy} if $\Pr(\tau)\ge2^{-\eta L}$. An
oblivious machine has a single trace, of probability~$1$.

If a run appends exactly $L$ bits before step $T$, let $a(r)$ be the step at which output bit $r$ is appended. For an input
position $q$ the \emph{root} of $q$ is
\[
  \Out(q)=a(\tpos(q)),
\]
the step at which the output copy of input bit $q$ is written, and $\Out(V)=\{\Out(q):q\in V\}$ for $V\subseteq\Range L$. The
map $\Out$ is injective, because append steps are distinct and $\tpos$ is a bijection
(\lean{appendTime_tposOf_injective}). Results quantified over sets $V_0$ of input positions are therefore statements about the
output steps of the corresponding output sets $\tpos(V_0)$.

\subsection{Revisits, excursions and crossings}

Fix a work tape, with head positions $y(t)$, and a horizon $T$.

\begin{definition}[Revisits]\label{def:hop}
\formal{EnvelopeCrux.lean}{BoundedNesting.IsHop, HasHopExcursionGTAt, HopsCross}
A \emph{revisit} (in Lean, a hop) is a pair of steps $p<u<T$ with $p+2\le u$, $y(u)=y(p)$, and $y(t)\ne y(p)$ for all $p<t<u$:
two consecutive visits to one cell between which the head leaves it. Its \emph{excursion} is
$\max_{p<t<u}\abs{y(t)-y(p)}$, the farthest the head travels from the cell before returning. Its \emph{earlier end} is~$p$.
Two revisits $(p,u)$ and $(p',u')$ of the same tape \emph{cross} if $p<p'<u<u'$ or $p'<p<u'<u$.
\end{definition}

A revisit lies on one side of its cell: the head stays strictly above it, or strictly below it, in between. Two crossing
revisits lie on opposite sides (\cref{lem:zigzag}(a); compare~\cite[\S2]{EighthPaper2026}). Every step is the earlier end of at most one revisit on each tape.

\begin{definition}[Crossing steps]\label{def:crossing}
\formal{CrossingBridge.lean}{crossingLowerEtaSteps}
For a trace with horizon $T$ and an integer $H\ge1$, let $\CL_H$ be the set of steps $v<T$ such that, on some work tape, $v$ is
the earlier end of a revisit of excursion greater than $2H$ that crosses another revisit of excursion greater than $2H$ on the same
tape.
\end{definition}

In Lean $\CL_H$ is \lean{crossingLowerEtaSteps}~$\tau\,T\,H\,1\,T$. The set shrinks as $H$ grows
(\lean{crossingLowerEtaSteps_anti_H}), so a lower bound on the cost of separating from $\CL_H$ is stronger for larger~$H$.

To see what $\CL_H$ measures, consider a head that sweeps right across a stretch of tape, turns, sweeps back, turns again and
sweeps right. A cell visited on the first sweep and revisited on the second is the earlier end of a revisit above it; a cell
visited on the second sweep and revisited on the third is the earlier end of a revisit below it; and when the second turn
overlaps the first stretch these revisits cross. This is how a tape is used as a queue. A tape used as a stack, written while the
head moves one way and read back while it returns, makes only nested revisits. \Cref{sec:skeletons} turns this picture into
theorems.

\subsection{The dependency DAG, cuts and cone relations}

\begin{definition}[Dependency DAG]\label{def:dag}
\formal{MergeDepth/ColdGenie.lean}{MergeDepth.DepStep, MergeDepth.IsDepPath}
The \emph{dependency DAG} of a trace with horizon $T$ has the vertices $0,\dots,T-1$ and two kinds of edges: the \emph{state edge}
$t\to t-1$ for $1\le t<T$, and the \emph{revisit edge} $u\to p$ for every revisit $(p,u)$ of every work tape. A \emph{dependency
path} is a sequence $t_0\to t_1\to\dots\to t_m$ of edges.
\end{definition}

Every edge decreases the step, so the DAG is acyclic. The value written at step $t$ is a function of the values at its
out-neighbours and of the input bit read at~$t$: the state edge carries the state, and the revisit edges carry the contents of the
scanned cells. In Lean the edge relation is \lean{MergeDepth.DepStep} and paths are \lean{MergeDepth.IsDepPath} (from the eleventh
paper's library~\cite{MergeDepth2026}).

\begin{definition}[Cuts]\label{def:cut}
\formal{OneFreeHeadSmallExcursion.lean}{IsMengerVertexCut}
For sets $S$ and $D$ of steps, a set $B$ of steps \emph{separates $S$ from $D$} if every dependency path that starts in $S$ and
ends in $D$ contains a vertex of $B$, its endpoints included. In Lean this is \lean{IsMengerVertexCut}.
\end{definition}

By Menger's theorem the least size of a set separating $S$ from $D$ equals the largest number of vertex-disjoint dependency paths
from $S$ to $D$. We use Menger's theorem only to explain statements; the theorems and their formal proofs bound the size of
separating sets directly.

\begin{definition}[Cones, balls and acyclic sets]\label{def:cone}
\formal{EnvelopeCrux.lean}{BoundedStackDepth.depCone, BoundedStackDepth.coneSystem, ballAvoidSystem, HalfRate.IsAcyclicIn, DirectedRecovery.IsIndependent}
Let $\Gamma$ be a set of steps. For a step $t$, the \emph{cone} of $t$ avoiding $\Gamma$ is the set of input positions $p$ such
that some dependency path $t\to t_1\to\dots\to t_m$ with $t_1,\dots,t_m\notin\Gamma$ ends at a step $t_m$ (or at $t$ itself, if
$m=0$) at which the input head reads position $p$. The root $t$ itself need not avoid~$\Gamma$. The \emph{cone relation}
assigns to each input position $q$ the cone of its root $\Out(q)$ (in Lean, \lean{BoundedStackDepth.coneSystem} evaluated at $\tpos(q)$). The
\emph{ball relation of radius $D$ avoiding $B$} is defined in the same way using only paths with at most $D$ edges after the
root (\lean{ballAvoidSystem}). A set $U$ of input positions is \emph{acyclic} in a relation $N$ if every nonempty $S\subseteq U$
has an element $v$ with $N(v)\cap S=\emptyset$ (\lean{HalfRate.IsAcyclicIn}). It is \emph{independent} in $N$ if
$N(v)\cap U=\emptyset$ for every $v\in U$ (\lean{DirectedRecovery.IsIndependent}); an independent set is acyclic.
\end{definition}

An acyclic set is a set of receivers none of which can see, in its cone, the bits wanted by the receivers that come after it in
a suitable order. In index coding, a message that lets every receiver of an acyclic set recover its bit must carry at least as
many bits as the set has elements~\cite{BBJK2011}. For machines this takes the following form: what the cut steps write, together
with what a small \emph{genie} reveals, is a message from which each output of the acyclic set is computed using only its cone.
The fourteenth paper's genie (\cref{lem:genie} below) bounds the depth of the cones at small cost, and the eleventh paper's
index-code lemma turns an oblivious transposer into such a message~\cite{MergeDepth2026}. For the cone relation itself, which has
no depth bound, no genie is needed (\cref{lem:acyclic-cost}(b)).

\section{The certificate method}\label{sec:certificates}

Every exclusion in this paper has the same shape: find, in the cone relation of a run, a linear-size acyclic set, possibly after
revealing a cheap set of steps, and conclude that the run cannot belong to a linear-time transposer. This section states the
method and the tools used to find acyclic sets.

\subsection{Genies and index codes}

\begin{lemma}[Genies and index codes]\label{lem:genie}
\cited{\cite[Corollary~O]{BoundedStack2026}, Lean-verified there (\lean{BoundedStackDepth.genie_cost},
\lean{BoundedStackDepth.genie_of_trace}, \lean{BoundedStackDepth.exists_indexCode}).}
\begin{enumerate}[label=(\alph*)]
\item For every $k\ge1$, $\beta\ge1$, $C\ge1$, $\eps>0$ and every trace with $k$ work tapes and $T\le C(L+1)$, there is a set $\Gamma$
of steps with $\beta\abs\Gamma\le\eps L$ such that every dependency path whose vertices after the first avoid $\Gamma$ has at most
$\mathrm{genieDepth}(k,\beta,C,\eps)$ vertices. The depth bound depends only on $k$, $\beta$, $C$ and~$\eps$.
\item If $M$ is oblivious at length $L$ with common trace $\tau$ and transposes, then for every set $\Gamma$ of steps the values
written at the steps of $\Gamma$ form an index code of $\wval\abs\Gamma$ bits for the cone relation avoiding $\Gamma$: two inputs
with the same values on $\Gamma$ that agree on the cone of $\Out(q)$ agree on the output bit at $\Out(q)$.
\end{enumerate}
\end{lemma}

The genie comes from the fourteenth paper's solution of the trajectory cut problem~\cite[Theorem~T]{BoundedStack2026}; its depth
$\mathrm{genieDepth}(k,\beta,C,\eps)$ is explicit, with $\log\log\mathrm{genieDepth}=O(k^2\log(k\beta C/\eps))$~\cite[Remark~9.4 and
\S10]{BoundedStack2026}.

\begin{lemma}[Acyclic sets bound codes]\label{lem:acyclic-code}
\formal{DichromaticCrux.lean}{storageCode_card_le_of_isAcyclicIn, indexCode_bits_ge_of_isAcyclicIn}
Let $N$ be a relation on a finite set $V$ and $A\subseteq V$ acyclic in~$N$. Every storage code for $N$ over an alphabet $\Sigma$ has
at most $\abs\Sigma^{\abs V-\abs A}$ words, and every index code for $N$ and the demand $\tpos$ has at least $\abs A$ bits, provided
$A$ is acyclic for $q\mapsto N(\tpos(q))$.
\end{lemma}

\begin{outline}
On a storage code a word is determined by its letters outside $A$: recover the letters in $A$ in an order in which each letter's
neighbourhood lies outside the not-yet-recovered part (\lean{acyclicIn_outside_injOn}). The fibres of an index code are storage
codes for the demand-shifted relation, and pigeonhole over the messages gives the bound.
\end{outline}

\subsection{Cruxes}

\begin{definition}[Cruxes]\label{def:crux}
\formal{MachineCrux.lean, EnvelopeCrux.lean}{SubpopCutCruxAt, MachineCutCruxAt, MachineCruxAt, subpopCutCruxAt_one_iff, machineCutCruxAt_zero_iff}
Let $M$ be a machine, $C\ge1$, $D\ge0$, $0<\alpha_0\le1$ and $\rho,\delta\ge0$. $M$ has \emph{subpopulation cut-crux rate $\rho$ at
density $\alpha_0$, cost $\delta$ and depth $D$} if for every $\eta>0$ and all large $d$, whenever $M$ is oblivious at length $L$
with halting time $T\le C(L+1)$ and exactly $L$ appends, there are a set $V_0$ of at least $\alpha_0L$ input positions and a set $B$
of steps with $\wval\abs B\le\delta\abs{V_0}$ such that every index code for the cone relation avoiding $\Gamma\cup B$, for any
$\Gamma$ that cuts all dependency paths to at most $D+1$ vertices, has at least $(\rho-\eta)\abs{V_0}$ bits (\lean{SubpopCutCruxAt}).
Two cases have their own names.
\begin{enumerate}[label=(\alph*)]
\item At density $1$ and cost $0$, $M$ has \emph{crux rate $\rho$ at depth $D$}. Then $V_0=[0,L)$ and $B=\emptyset$, since $\wval\ge2$:
every index code for the cone relation avoiding $\Gamma$ has at least $(\rho-\eta)L$ bits (\lean{MachineCruxAt},
\lean{machineCutCruxAt_zero_iff}).
\item At density $1$, $M$ has \emph{cut-crux rate $\rho$ at cost $\delta$ and depth $D$}. Then $V_0=[0,L)$: there is a set $B$ with
$\wval\abs B\le\delta L$ such that every index code for the cone relation avoiding $\Gamma\cup B$ has at least $(\rho-\eta)L$ bits
(\lean{MachineCutCruxAt}, \lean{subpopCutCruxAt_one_iff}).
\end{enumerate}
\end{definition}

\begin{theorem}[Cruxes exclude]\label{thm:crux-excludes}
\formal{EnvelopeCrux.lean, MachineCrux.lean}{subpop_cutcrux_excludes, machine_cutcrux_excludes, machine_crux_excludes}
Let $C\ge1$ and $\eps>0$, and put $D_\eps=\mathrm{genieDepth}(k,\wval,C,\eps)$. If $M$ has subpopulation cut-crux rate $\rho$ at
density $\alpha_0$, cost $\delta$ and depth $D_\eps$ with $\eps<\alpha_0(\rho-\delta)$, then for all large $d$ there is no $T$ with
$\mathrm{ObliviousLinTime}(M,d,C,T)$. In particular a cut-crux rate $\rho$ at cost $\delta$ with $\eps+\delta<\rho$, and a crux rate
$\rho>\eps$, at depth $D_\eps$ exclude.
\end{theorem}

\begin{outline}
Suppose $M$ transposes obliviously in time $C(L+1)$. Reveal the genie $\Gamma$ of \cref{lem:genie}(a), of cost $\eps L$ and depth
$D_\eps$, together with~$B$. By \cref{lem:genie}(b) the values on $\Gamma\cup B$ form an index code of at most $\eps L+\delta\abs{V_0}$
bits, so $(\rho-\delta-\eta)\abs{V_0}\le\eps L$, which contradicts $\abs{V_0}\ge\alpha_0L$ for small $\eta$. The two special cases
are $\alpha_0=1$, where $\eps<\alpha_0(\rho-\delta)$ reads $\eps+\delta<\rho$, and in addition $\delta=0$, where it reads $\eps<\rho$;
their direct forms are \lean{machine_cutcrux_excludes} and \lean{machine_crux_excludes}.
\end{outline}

Acyclic sets give cruxes. If, for every $\eta>0$ and all large $d$, every cone relation of depth $D$ in \cref{def:crux}(a) contains
an acyclic set of at least $(\rho-\eta)L$ input positions, then $M$ has crux rate $\rho$ at depth~$D$, by \cref{lem:acyclic-code}
(\lean{machineCruxAt_of_isAcyclicIn}, \lean{machineCrux_of_isAcyclicIn}, \lean{oblivious_excludes_of_isAcyclicIn}). The same holds for ball relations, which contain the
cone relations of depth~$D$ (\lean{machineCutCruxAt_of_ballAvoid_isAcyclicIn}).

Part~(b) of the next lemma is the form used by Theorems~X and~S: at a single size, revealing a set $B$ that leaves an acyclic set $U$ in
the cone relation costs at least $\abs U$ bits. Part~(a), used by the first forms, does the same for the ball relation of a bounded
radius, up to the genie.

\begin{lemma}[Acyclic sets cost cuts]\label{lem:acyclic-cost}
\formal{EnvelopeCrux.lean, TightConstants.lean}{obliviousLinTime_separator_card_ge_of_isAcyclicIn, obliviousLinTime_coneSystem_isAcyclicIn_card_le, obliviousLinTime_ballSeparator_card_ge}
Let $\mathrm{ObliviousLinTime}(M,d,C,T)$ with $d\ge1$, $C\ge1$ and $\eps>0$, put $D_\eps=\mathrm{genieDepth}(k,\wval,C,\eps)$, and let
$B$ be a set of steps.
\begin{enumerate}[label=(\alph*)]
\item For every set $U$ of input positions that is acyclic in the ball relation of radius $D_\eps$ avoiding $B$,
$\wval\abs B\ge\abs U-\eps L$.
\item For every set $U$ of input positions that is acyclic in the cone relation avoiding $B$, $\wval\abs B\ge\abs U$; this part needs
only $\mathrm{ObliviousLinTime}(M,d,C,T)$.
\end{enumerate}
\end{lemma}

\begin{outline}
(a) Reveal the genie of cost $\eps L$ together with $B$. The cone relation avoiding both is contained in the ball relation
(\lean{coneSystem_subset_ballAvoidSystem}), so $U$ is acyclic in it, and \cref{lem:genie}(b) with \cref{lem:acyclic-code} gives
$\wval(\abs\Gamma+\abs B)\ge\abs U$. (b) The values on $B$ alone already form an index code for the cone relation avoiding $B$,
which has no depth bound, and \cref{lem:acyclic-code} applies with no genie.
\end{outline}

\subsection{Heavy traces of arbitrary machines}

For a machine that is not oblivious, the index code of \cref{lem:genie}(b) is not available, but counting the fibre of a single
trace replaces it.

\begin{lemma}[Frozen-trace deficit]\label{lem:frozen}
\formal{SimulatorBridge.lean}{correctSubset_isAcyclicIn_coneSystem_deficit_le_logb, correctSubset_isAcyclicIn_coneSystem_deficit_le, frozenTrace_isAcyclicIn_coneSystem_deficit_le, frozenTrace_isAcyclicIn_ballAvoid_separator_card_ge}
Let $M$ be any machine, $d\ge1$, $\tau$ a trace of $M$ with halting time $T$, and $S$ a nonempty set of inputs that follow $\tau$ and
that $M$ transposes within time $c$. Then $T\le2cL$, $\tau$ appends exactly $L$ bits, and for every set $\Gamma$ of steps and every set
$U$ acyclic in the cone relation of $\tau$ avoiding $\Gamma$,
\[
  \abs U-\wval\abs\Gamma\ \le\ L-\log_2\abs S .
\]
In particular $\abs U-\wval\abs\Gamma\le\eta L$ if $\abs S\ge2^{(1-\eta)L}$; for instance, if $M$ transposes at size $d$ within time
$c$ and $\tau$ is $\eta$-heavy, $S$ can be the set of all inputs that follow $\tau$.
\end{lemma}

\begin{outline}
All inputs of $S$ share the dependency DAG of $\tau$, and any one of them gives the facts about $T$. Split $S$ according to the
$2^{\wval\abs\Gamma}$ possible values on $\Gamma$. Since the inputs of $S$ are transposed, the bit appended at $\Out(q)$ is input
bit $q$, so each part is a storage code for the cone relation and has at most $2^{L-\abs U}$ elements by \cref{lem:acyclic-code}.
Hence $\abs S\le2^{\wval\abs\Gamma+L-\abs U}$. Correctness is used only on~$S$.
\end{outline}

\subsection{Tools for finding acyclic sets}

Acyclic sets in cone relations are found by potentials. The following statements hold for any relation $N$ on a finite set; an arc
$q\to p$ means $p\in N(q)$.

\begin{lemma}[Envelopes and windows]\label{lem:window}
\formal{EnvelopeCrux.lean}{acyclicIn_of_lt_on_nbhd, upperEnvelope_isAcyclicIn, oneSidedEnvelope_isAcyclicIn, exists_isAcyclicIn_of_oneSidedEnvelope_maxDeg, window_isAcyclicIn_of_noFarArcs, exists_isAcyclicIn_of_window_maxDeg}
\begin{enumerate}[label=(\alph*)]
\item If a potential strictly decreases along every arc inside $S$, then $S$ is acyclic.
\item Let $\varphi,\psi$ be real potentials and $w$ a real number. The \emph{window} $\{q\in V_0:q\notin N(q),\ \varphi(q)+w<\psi(q)\}$
is acyclic if every arc $q\to p$ inside it is \emph{near}, $\psi(p)\le\varphi(q)+w$. If instead every vertex of the window is
incident to at most $\beta$ far arcs inside it, the window contains an acyclic set of at least a $1/(\beta+1)$ fraction of its
elements.
\end{enumerate}
\end{lemma}

\begin{outline}
(a) Inside $S$ an element of least potential has no out-arc into $S$. (b) A near arc $q\to p$ inside the window has
$\varphi(p)<\psi(p)-w\le\varphi(q)$, so $\varphi$ strictly decreases along it and (a) applies. With far arcs, remove the far neighbours
of the vertices one at a time; a greedy choice keeps a $1/(\beta+1)$ fraction. A two-sided bound $\abs{\psi(p)-\varphi(q)}\le w$ is
near at margin $2w$; with two-sided bounds the statements are \lean{window_isAcyclicIn_of_noFarArcs} and
\lean{exists_isAcyclicIn_of_window_maxDeg}. The upper-envelope form of (a), that $\{q\in V_0:\Psi(q)<\varphi(q)\}$ is acyclic when
$\varphi(r)\le\Psi(p)$ for every arc $r\to p$ inside $V_0$, is \lean{upperEnvelope_isAcyclicIn}.
\end{outline}

\begin{lemma}[Cut sets and block ranking]\label{lem:blockcut}
\formal{EnvelopeCrux.lean}{cutSet_isAcyclicIn, exists_isAcyclicIn_card_ge_of_blockCut, exists_isAcyclicIn_card_ge_of_uniformBlockCut, exists_isAcyclicIn_of_shiftBlockCutFamily_gen}
Assign to each $q\in V_0$ two disjoint sets $A(q),E(q)$ of blocks in $\{0,\dots,B-1\}$, and suppose every arc $q\to p$ inside $V_0$ has
$A(q)\cap E(p)\ne\emptyset$. For a ranking $\sigma$ of the blocks, the set of $q$ with $\sigma(A(q))<\sigma(E(q))$ (every block of $A(q)$
ranked below every block of $E(q)$) is acyclic, and some ranking makes it contain at least
$\sum_q\binom{\abs{A(q)}+\abs{E(q)}}{\abs{A(q)}}^{-1}$ elements, hence at least $\abs{V_0}/\binom{a+e}{a}$ elements when
$\abs{A(q)}\le a$ and $\abs{E(q)}\le e$ throughout.
\end{lemma}

\begin{outline}
Along an arc $q\to p$ inside the set, a block of $A(q)\cap E(p)$ shows that the largest rank in $A(p)$ is below the smallest rank in $E(p)$,
hence below that block's rank, hence at most the largest rank in $A(q)$; so the largest rank of $A$ strictly decreases along arcs, and
\cref{lem:window}(a) applies. For disjoint $A(q),E(q)$, exactly a fraction $\binom{\abs{A(q)}+\abs{E(q)}}{\abs{A(q)}}^{-1}$ of all rankings
put $A(q)$ below $E(q)$ (\lean{LeanMachine.BlockRanking.choose_mul_card_goodPerms_eq_factorial}), so the average size of the set over all rankings is the stated
sum. For a family of $m$ block assignments, of which $s$ elements in total exceed the sizes $a,e$ and at most $\kappa$ per assignment
have $A(q)\cap E(q)\ne\emptyset$, the assignment with the fewest exceptions gives an acyclic set of at least
$(\abs{V_0}-s/m-\kappa)/\binom{a+e}{a}$ elements (\lean{exists_isAcyclicIn_of_shiftBlockCutFamily_gen}).
\end{outline}

\begin{lemma}[Step-invariant windows]\label{lem:stepinv}
\formal{EnvelopeCrux.lean}{isAcyclicIn_ballAvoid_window_of_stepInvariant, isAcyclicIn_ballAvoid_window_of_constStepInvariant}
Let a trace, a set $B$ of steps, a depth $D$ and a set $V_0$ of input positions be given, with regions $R(0),\dots,R(D)$ of steps, a
\emph{sink} region, a step potential $\mu$, potentials $\varphi,\psi$ on input positions, reals $c_0,c_1$ and $\Delta\ge0$. Suppose:
the roots of $V_0$ lie in $R(D)$ with $\mu(\Out(q))\le\varphi(q)+c_0$; every edge from the sink that avoids $B$ stays in the sink,
and the sink reads no position of $V_0$; every edge $u\to v$ from $R(r+1)$ that avoids $B$ enters the sink or enters $R(r)$ with
$\mu(v)\le\mu(u)+\Delta$; and a read of $p\in V_0$ at a step $u\in R(r)$ has $\psi(p)\le\mu(u)+c_1$. Then
$\{q\in V_0:\varphi(q)+c_0+D\Delta+c_1<\psi(q)\}$ is acyclic in the ball relation of radius $D$ avoiding~$B$.
\end{lemma}

\begin{outline}
By induction along a dependency path from a root with at most $D$ edges that avoids $B$ after the root, its end lies in the sink or in
the region $R(D-\ell)$ after $\ell$ edges, with $\mu$ raised by at most $\ell\Delta\le D\Delta$
(\lean{depPath_oneSidedStepInvariant_with_sink}). Since the sink reads nothing in $V_0$, every arc $q\to p$ inside $V_0$ has
$\psi(p)\le\varphi(q)+c_0+D\Delta+c_1$: it is near, and \cref{lem:window}(b) applies. The strict inequality excludes loops, so no
sign condition on $c_0,c_1$ is needed. The condition $\Delta\ge0$ is needed: for $D=1$ and $\Delta=-1$, a root that reads its own
position is a loop inside the set.
\end{outline}

\Cref{lem:window,lem:blockcut,lem:stepinv} are statements about traces only. They are the tools behind the exclusions of the
schedule families in \cref{app:catalogue}. The first form of \cref{thm:pullback} (\lean{timePullback_blockCut_mais}) was a far-reaching use of \cref{lem:blockcut};
the present form replaces the uniform block cut by a single interval on which the block sets are constant.

\section{Two long-range heads}\label{sec:heads}

The first universal consequence of the certificate method concerns how far the heads must travel. It has polynomial constants
and holds for every trace that carries a fixed exponentially small fraction of the inputs.

For reals $c,\beta\ge1$ and an integer $k\ge1$ put
\[
  N=\lfloor160\beta c^2\rfloor,\qquad \eps_*=\frac1{8(N+1)},\qquad \eps_{\mathrm{sm}}=\frac{\eps_*}{4\beta(N+2)\max(1,2k)},
\]
\[
  d_{\mathrm{H}}=\max\Bigl(\lceil80\beta c\rceil,\ \Bigl\lceil\frac{20(N+2)}{\eps_*}\Bigr\rceil,\
  \Bigl\lceil\frac{16\beta(N+2)k}{\eps_*}\Bigr\rceil,\ 10\Bigr)
\]
(\lean{slabWallCap}, \lean{slabAutonomyEps}, \lean{slabSmallExcursion}, \lean{slabDimThreshold}). For $c=1$, $\beta=2$ and $k=2$:
$\eps_*\approx3.9\cdot10^{-4}$, $\eps_{\mathrm{sm}}\approx3.8\cdot10^{-8}$ and $d_{\mathrm{H}}\approx5.3\cdot10^7$.

\begin{theorem}[Theorem~H, two long-range heads]\label{thm:H}
\formal{OneFreeHeadSmallExcursion.lean}{oneFreeHead_smallExcursion_excludes, twoMacroscopicHeads_of_linearTimeTransposer, twoMacroscopicHeads_of_oblivious_linearTimeTransposer, offCarrier_vertexExpansion_of_linearTimeTransposer, offCarrier_vertexExpansion_of_obliviousLinTime}
Let $c,\beta\ge1$ and $k\ge1$. Let $M$ have at most $k$ work tapes and value width at most $\beta$, let $d\ge d_{\mathrm{H}}$, and suppose $M$
transposes at size $d$ within time~$c$. Let $\tau$ be an $\eta$-heavy trace of $M$ with halting time $T$; for an oblivious $M$, take its
common trace and $\eta=0$.
\begin{enumerate}[label=(\alph*)]
\item \emph{Slab autonomy.} If $\eta<\eps_*$, there is no set of at most one work head such that every work head outside it has all its
revisits of excursion at most $\lfloor\eps_{\mathrm{sm}}L\rfloor$ on~$\tau$. The input head and the excepted head are unrestricted.
\item \emph{Two heads.} If $\eta<\eps_*$, two distinct work heads each make, on $\tau$, a revisit of excursion greater than
$\eps_{\mathrm{sm}}L$.
\item \emph{Off-carrier expansion.} For every set $S$ of at most one work head and every $H\le\lfloor\eps_{\mathrm{sm}}L\rfloor$, there is a
set $V_*$ of input positions with $3\eps_*L\le\abs{V_*}\le L/2$ such that every set $B$ of steps separating $\Out(V_*)$ from the earlier
ends of revisits of excursion greater than $H$ on the work heads outside $S$ satisfies
\[
  (\eps_*-\eta)L\ \le\ \beta\abs B+\abs{B\cap\Out(V_*)} .
\]
\end{enumerate}
\end{theorem}

Part (b) is part (a) applied to the set of heads that make long revisits. In (c) the set $V_*$ depends on the trace and on $S$; unlike
Theorem~X, (c) does not hold for every large output set, and its targets are all long revisits off one head, not only crossing
ones. It has no hypothesis on $\eta$, but is vacuous for $\eta\ge\eps_*$.

\begin{outline}
\emph{Frozen-trace chunks.} If every dependency path that starts in a set $C$ of steps and avoids $\Gamma$ after its first vertex
stays in $C$, then the input positions whose output copies are appended in $C$, minus those read in $C$, form an independent set
$U_*$, and counting the fibre as in \cref{lem:frozen} gives $\abs{U_*}-\beta\abs\Gamma\le\eta L$ (\lean{frozenTrace_chunk_deficit_le}).
\emph{Slabs.} Cut the tape of the unrestricted head into slabs of width about $L/(80\beta c)$; some offset makes it cross slab walls at
most $N$ times (\lean{exists_offset_slabTransitions_le}). \emph{A chunk.} Some time interval inside one slab epoch, of at most $L/2$
steps, appends at least $3\eps_*L$ output bits (\lean{exists_slab_chunk_dichotomy}). \emph{Trapping.} A genie made of the slab
transitions, the last visits of the free head before it leaves the slab, and the earlier ends of the other heads' revisits that span
the chunk boundaries traps every path from the chunk (\lean{depPath_trapped_in_chunk}); the last part is small because those
revisits are short, and a short revisit can span a given step on at most $2H+1$ cells (\lean{card_smallCrossingLowers_le}).
\emph{Counting.} The chunk has at most $N+1$ runs, and every output row meets every input row in $b$ bits, so the independent set
keeps about half the outputs of the chunk (\lean{outIn_card_le_two_indep_add_six_runs}). The resulting deficit $\eps_*L$ exceeds
$\eta L$, a contradiction. For (c), keep in the genie only the short revisits, which are always few, and add the separating set $B$:
paths that avoid both stay in the chunk, and the same count gives the inequality.
\end{outline}

\begin{corollary}[Sublinear excursions, family heights, one tape]\label{cor:H}
\formal{OneFreeHeadSmallExcursion.lean}{oneFreeHead_sublinearExcursion_excludes, oneFreeHead_sublinearExcursion_oblivious_excludes, oneFreeHead_sublinearFamHeight_excludes, atMostOneWorkTape_lowAdaptivity_excludes, atMostOneWorkTape_oblivious_excludes}
Let $c,\beta\ge1$, $k\ge1$ and $H(d)=o(L)$. For all large $d$, no machine with at most $k$ work tapes and value width at most $\beta$
that transposes at size $d$ within time $c$ has an $\eta$-heavy trace, $\eta<\eps_*$, on which all but one work head make only revisits
of excursion at most $H(d)$; the same holds with the excursion replaced by the height of every family of revisits that share a side
and a turning point. In particular no such machine has at most one work tape.
\end{corollary}

\begin{outline}[Outline of the formal proof of \cref{cor:H}]
For $d$ large, $H(d)\le\eps_{\mathrm{sm}}L$, so \cref{thm:H}(a) applies; the height of a family bounds the excursions of its revisits
(\lean{hopExcursionLEOn_of_famd_le}). With at most one work tape, take the excepted set to contain every work head.
\end{outline}

Superlinear lower bounds for one work tape are classical in the one-tape models~\cite{DMS1991,DietzfelbingerMaass1993}, and in the present
model the fourth paper's analysis of macroscopic limits already excludes one work tape~\cite[Corollary~3.18(c)]{FourthPaper2026}; we
include the case only because it falls out. The fourth paper also showed, by a written proof, that in every macroscopic limit of a
linear-time transposer at least two work-head limits are non-monotone~\cite[Corollaries~3.18 and~3.24]{FourthPaper2026}.
Theorem~H(b) is a quantitative form of this at a single size, with explicit polynomial constants, on every $\eta$-heavy trace with
$\eta<\eps_*$. Theorem~H(a) excludes one free head even when the other heads make excursions up to a small fixed fraction of $L$, which the
limit statement does not exclude, and part~(c) is new. Every oblivious linear-time transposer therefore has at least two work tapes on
which information is stored and retrieved across distances that are a fixed fraction of~$L$.

\section{Macroscopic crossings}\label{sec:crossings}

This section proves the dichotomy for arbitrary linear-time traces (Theorem~D) and its consequence for oblivious transposers
(Theorem~X). Both rest on one implication, proved for every trace without any hypothesis on what the machine computes: if a
small set of steps cuts a large set of outputs off from the crossing steps $\CL_H$, then the remaining cone relation contains a
linear-size acyclic set. \Cref{sec:certificates} has shown that such an acyclic set is incompatible with transposition.

\subsection{Constants}\label{sec:constants}

Fix $k\ge1$, reals $C\ge1$ and $\alpha>0$, and a depth $D\ge0$. Put
\[
  m_0(k)=(2k+1)(k+1)^{2k},\qquad K_{\mathrm{row}}=\lceil24C/\alpha\rceil,\qquad B_{\mathrm{time}}=\lceil8m_0C/\alpha\rceil,
\]
\[
  \rho_0(k,C,\alpha)=\frac{1}{3\binom{m_0+2K_{\mathrm{row}}}{m_0}},\qquad
  N_{\mathrm{thr}}=(2k+1)+m_0(2k+1)(B_{\mathrm{time}}+1),\qquad
  E_0(k,D,C,\alpha)=(k+2)^D N_{\mathrm{thr}}(5k+1),
\]
\[
  \theta_{\mathrm{E}}(k,D,C,\alpha)=\frac{\alpha}{12\,E_0(k,D,C,\alpha)},\qquad
  d_{\mathrm{pb}}(k,C,\alpha)=\max\Bigl(K_{\mathrm{row}}B_{\mathrm{time}},\Bigl\lceil\frac{12m_0(2B_{\mathrm{time}}+3)}{\alpha}\Bigr\rceil\Bigr).
\]
Here $m_0(k)$ counts the \emph{phase-exit chains} of the proof below; $m_0=12$, $405$ and $28672$ for $k=1,2,3$. The constant
$\rho_0$ does not depend on $D$ or on the alphabet. For a machine $M$ with $k$ work tapes and value width $\wval$, and $\alpha>0$,
put
\[
  \cE(M,C,\alpha)=\frac{\rho_0(k,C,\alpha)}{2\wval},\qquad
  D_{\mathrm{E}}(M,C,\alpha)=\mathrm{genieDepth}\bigl(k,\wval,C,\alpha\rho_0(k,C,\alpha)/4\bigr),
\]
\[
  \thetaM(M,C,\alpha)=\theta_{\mathrm{E}}\bigl(k,D_{\mathrm{E}}(M,C,\alpha/2),C,\alpha/2\bigr)
  =\frac{\alpha}{24\,E_0\bigl(k,D_{\mathrm{E}}(M,C,\alpha/2),C,\alpha/2\bigr)},
\]
where $\mathrm{genieDepth}$ is the depth bound of the fourteenth paper's genie (\cref{lem:genie}). In Lean these are
\lean{conjectureE_m0}, \lean{conjectureE_bridge_rho0} (through \lean{timePullbackRhoMin}), \lean{conjectureE_Nthresh},
\lean{conjectureE_E0}, \lean{conjectureE_theta}, \lean{timePullback_gen_d0}, \lean{conjectureE_cE}, \lean{conjectureE_DE} and
\lean{conjectureE_thetaMacro}.

These constants are positive and astronomically small. For $k=2$, $C=1$: $1/\rho_0(2,1,1)\approx10^{65.8}$ and
$1/\rho_0(2,1,\tfrac12)\approx10^{105.4}$; for $k=1$ and $k=3$ at $C=\alpha=1$ the values are $10^{12.6}$ and $10^{153.4}$. The
threshold $\thetaM$ is far smaller, because $E_0$ grows like $(k+2)^{D}$ and the genie depth $D_{\mathrm{E}}$ is itself at least
$32\wval C/(\alpha\rho_0)$: for $k=2$ and $C=\alpha=1$, $\thetaM\le4^{-10^{107}}$. Nothing in the problem suggests that the true
values are this small (\cref{rem:numerics-crossings}).

\subsection{The time-pullback bound}

The first ingredient turns a structure in time into an acyclic set. Split $\Range T$ into $B_{\mathrm{time}}$ consecutive
\emph{time blocks} $T_0,\dots,T_{B_{\mathrm{time}}-1}$ of length $\lfloor T/B_{\mathrm{time}}\rfloor$, the last one extended to end
at~$T$.

\begin{theorem}[Time pullback]\label{thm:pullback}
\formal{MacroscopicThreshold.lean}{theorem_51_2_general_zbad_mais}
Let $\tau$ be a trace with $k$ work tapes that appends exactly $L=d^2b$ bits before step $T\le C(L+1)$, with $d,b,m_0\ge1$, $C\ge1$ and
$\alpha>0$, and let $d\ge d_{\mathrm{pb}}$ computed with this $m_0$. Let $V_0\subseteq\Range L$ with $\abs{V_0}\ge\alpha L$, let
$Z\subseteq\Range L$ with $\abs Z\le\abs{V_0}/12$, let $B$ be a set of steps and $D\ge0$. Suppose there are at most $m_0$ indices
$\nu$, each with a nondecreasing sequence $\xi_\nu(0)\le\dots\le\xi_\nu(B_{\mathrm{time}})$ and a permutation $\pi_\nu$ of the
blocks, such that the following holds for every $q\in V_0\setminus Z$: whenever a dependency path from $\Out(q)$ with at most $D$
edges, whose vertices after the root avoid $B$, ends at a step of block $T_\beta$ at which the input head reads a position of $V_0$,
the root $\Out(q)$ lies in $[\xi_\nu(\pi_\nu(\beta)),\xi_\nu(\pi_\nu(\beta)+1))$ for some~$\nu$. Then the ball relation of radius
$D$ avoiding $B$ contains an acyclic set $U\subseteq V_0$ with
\[
  \abs U\ \ge\ \frac{\abs{V_0}}{3\binom{m_0+2K_{\mathrm{row}}}{m_0}} .
\]
\end{theorem}

\begin{outline}
For each $q$ let $A(q)$ be the set of blocks $\beta$ such that $\Out(q)$ lies in an interval assigned to $\beta$; since the
intervals of one $\nu$ are disjoint, $\abs{A(q)}\le m_0$. For an input position $p$ let $E(p)$ be the set of blocks in which $p$
is read. The hypothesis says that if $p$ is in the ball of $q$ then $A(q)$ and $E(p)$ meet. A kinematic count
(\lean{kinematic_goodInputRows_card_ge}) shows that all but at most $3T/(K_{\mathrm{row}}W)$ input rows are touched by the input head
in at most $K_{\mathrm{row}}$ runs, each shorter than a block, so $\abs{E(p)}\le2K_{\mathrm{row}}$ off a quarter of $V_0$. Pairs
with $A(q)\cap E(q)\ne\emptyset$, self-collisions, are bounded without any anti-concentration assumption, using the row and column
structure of $\tpos$ and the monotonicity of append steps: for a fixed $\nu$ each threshold splits at most one output row
(\lean{card_selfCollisions_singleFamily_le}). Off the bad sets the relation is contained in the \emph{block relation}
$A(q)\cap E(p)\ne\emptyset$ with $\abs A\le m_0$, $\abs E\le2K_{\mathrm{row}}$ and no self-loops, and a uniform block cut
(\lean{exists_isAcyclicIn_card_ge_of_uniformBlockCut}) extracts an acyclic set of the stated size.
\end{outline}

\Cref{thm:pullback} is purely about the trace. The constant $1/\binom{m_0+2K_{\mathrm{row}}}{m_0}$ is the price of the uniform
block cut over all possible block sets, and it is the source of the size of $\rho_0$.

\subsection{The bridge}

The second ingredient finds the time-pullback structure whenever $\CL_H$ can be cut off. We need a separation notion that
ignores the roots themselves.

\begin{definition}[Separation within depth $D$]\label{def:ballsep}
\formal{EnvelopeCrux.lean}{IsBallPathSeparator}
Let $V_0$ be a set of input positions and $S$, $B$ sets of steps. $B$ \emph{separates $V_0$ from $S$ within depth $D$} if every
dependency path from a root $\Out(q)$, $q\in V_0$, with at most $D$ edges and with all its vertices after the root outside $B$,
also has all its vertices after the root outside $S$ (\lean{IsBallPathSeparator}).
\end{definition}

If $B$ separates $\Out(V_0)$ from $S$ in the sense of \cref{def:cut}, then $B$ separates $V_0'=\{q\in V_0:\Out(q)\notin B\}$
from $S$ within every depth (\lean{isBallPathSeparator_of_isMengerVertexCut}).

\begin{theorem}[Bridge]\label{thm:bridge}
\formal{MacroscopicThreshold.lean}{conjectureE_bridge_exists_isAcyclicIn_crossingLowerGT2H, conjectureE_bridge_exists_isAcyclicIn_crossingLowerH}
Let $\tau$ be a trace with $k$ work tapes that appends exactly $L=d^2b$ bits before step $T\le C(L+1)$, with $C\ge1$ and
$\alpha>0$. Let $D\ge0$, let $1\le H\le\theta_{\mathrm{E}}(k,D,C,\alpha)\,L$, and let $d\ge d_{\mathrm{pb}}(k,C,\alpha)$. If $V_0$ has
at least $\alpha L$ elements and $B$ separates $V_0$ from $\CL_H$ within depth $D$, then the ball relation of radius $D$ avoiding
$B$ contains an acyclic set $U\subseteq V_0$ with $\abs U\ge\rho_0(k,C,\alpha)\abs{V_0}$.
\end{theorem}

\begin{outline}
\emph{Two cuts per tape.} On every unit-speed walk there are two steps $\vartheta_+,\vartheta_-$ such that every revisit of
excursion greater than $H$ that crosses no other such revisit straddles one of them: the upper ones straddle $\vartheta_+$, the
lower ones $\vartheta_-$ (\lean{exists_two_cuts_covering_nonCrossingLargeHops}). \emph{Dichotomy along paths.} Along a dependency
path from $V_0$ that avoids $B$ after the root, every edge is a state edge, a revisit edge of excursion at most $2H$, or a revisit
edge whose revisit crosses nothing large and therefore straddles one of the $2k$ cut steps; the separation hypothesis excludes
large crossing revisits (\lean{depPath_nonCrossing_dichotomy_of_isBallPathSeparator_crossingLowerH}). \emph{Phases and chains.}
The cut steps split time into at most $2k+1$ work phases; a path visits the phases along one of at most $m_0(k)$ phase-exit chains
(\lean{card_chainIdx_le}), and along a chain each passage between phases is a strictly antitone partial bijection between times
(\lean{crossingCut_nonCrossing_strictAnti}). Pulling the block boundaries back along each chain gives, for every chain, monotone
root intervals of both parities, which is the structure required by \cref{thm:pullback} with $m_0=m_0(k)$. \emph{Guards.} Small
revisits and the state edge are handled by a guard set of at most $(5k+1)H$ steps per threshold
(\lean{card_microHopStepGuard_le}); the roots that reach a guard within $D$ edges form the exceptional set $Z$, of size at most
$(k+2)^D N_{\mathrm{thr}}(5k+1)H=E_0H\le\alpha L/12$ by the choice of~$H$ (\lean{card_ancestors_le}). \Cref{thm:pullback} applies.
\end{outline}

The second Lean declaration is the same statement for banded crossing sets with an explicit clause on overshooting revisits;
for the unbanded set $\CL_H$ the clause holds for every trace (\lean{clauseE1H_uncapped}). The form with $H=W$ is
\lean{conjectureE_bridge_exists_isAcyclicIn_crossingLower}.

\subsection{Theorem D: a dichotomy for every linear-time trace}

\begin{theorem}[Theorem~D]\label{thm:D}
\formal{MacroscopicThreshold.lean}{linearTimeTrace_crossingLowerGT2H_menger_dichotomy, linearTimeTrace_crossingLowerGT2H_ballSep_dichotomy}
Let $\tau$ be any trace with $k$ work tapes that appends exactly $L=d^2b$ bits before step $T\le C(L+1)$, with $C\ge1$ and
$\alpha>0$. Let $D\ge0$, $1\le H\le\theta_{\mathrm{E}}(k,D,C,\alpha/2)\,L$, $d\ge d_{\mathrm{pb}}(k,C,\alpha/2)$, and $c\le\frac12$.
For every set $V_0$ of at least $\alpha L$ input positions, exactly one of the following holds:
\begin{enumerate}[label=(\roman*)]
\item some set $B$ of fewer than $c\abs{V_0}$ steps separates $\Out(V_0)$ from $\CL_H$, and the ball relation of radius $D$ avoiding
$B$ contains an acyclic set $U\subseteq V_0$ with $\abs U\ge\frac12\rho_0(k,C,\alpha/2)\abs{V_0}$;
\item every set separating $\Out(V_0)$ from $\CL_H$ has at least $c\abs{V_0}$ elements.
\end{enumerate}
\end{theorem}

\begin{outline}
If (ii) fails, take a separating set $B$ with $\abs B<c\abs{V_0}\le\abs{V_0}/2$. Since $\Out$ is injective, at least half of the
roots of $V_0$ lie outside $B$; the set $V_0'$ of these has at least $(\alpha/2)L$ elements, and $B$ separates it from $\CL_H$
within depth~$D$. \Cref{thm:bridge} at density $\alpha/2$ gives $U$.
\end{outline}

Theorem~D mentions neither a machine nor transposition. Its content is the implication from (i)'s first half to its second: a
cheap cut from the crossing steps leaves a large acyclic set. Read the other way, it says that a trace whose cone relations have
no large acyclic set after cheap cuts must reach $\CL_H$ with linear expansion from every large output set. The transposition
hypothesis enters only through the certificate method.

\subsection{Theorem X}

\begin{theorem}[Theorem~X, macroscopic crossings]\label{thm:X}
\formal{MacroscopicThreshold.lean}{conjectureE_obliviousLinTime_macroscopicCrossingLowers_menger_lb, conjectureE_obliviousLinTime_crossingLowerGT2H_menger_lb}
Let $M$ be a machine with $k$ work tapes and value width $\wval$, and let $C\ge1$ and $\alpha>0$. Put $\theta=\thetaM(M,C,\alpha)$ and
\[
  d_0=\max\Bigl(1,\ d_{\mathrm{pb}}(k,C,\alpha/2),\ \Bigl\lceil\frac{24E_0(k,D_{\mathrm{E}}(M,C,\alpha/2),C,\alpha/2)}{\alpha}\Bigr\rceil,\
  \Bigl\lceil\frac2\theta\Bigr\rceil\Bigr).
\]
Let $d\ge d_0$, suppose $\mathrm{ObliviousLinTime}(M,d,C,T)$, and let $H^*=\lfloor\theta L\rfloor$. Then $W\le H^*$ and
$\theta L\le2H^*$, and for every set $V_0$ of at least $\alpha L$ input positions, every set of steps separating $\Out(V_0)$ from
$\CL_{H^*}$ in the dependency DAG of the common trace has at least
\[
  \tfrac12\,\cE(M,C,\alpha/2)\,\abs{V_0}\ =\ \frac{\rho_0(k,C,\alpha/2)}{4\wval}\,\abs{V_0}
\]
elements. The same bound holds for $\CL_H$ at every $1\le H\le\theta L$ and every $d\ge\max(1,d_{\mathrm{pb}}(k,C,\alpha/2))$.
\end{theorem}

Every revisit counted in $\CL_{H^*}$ has excursion greater than $2H^*\ge\theta L$: these are crossings of macroscopic revisits.
When $\alpha\le1$, taking $V_0$ to be a set of $\lceil\alpha L\rceil$ input positions and the separating set to be $\CL_{H^*}$ itself gives the
following.

\begin{corollary}[Linearly many macroscopic crossings]\label{cor:many}
\formal{MacroscopicThreshold.lean}{conjectureE_obliviousLinTime_macroscopicCrossingLowers_menger_lb}
Under the hypotheses of \cref{thm:X}, if moreover $\alpha\le1$, the common trace has at least $\frac{\alpha\rho_0(k,C,\alpha/2)}{4\wval}L$ steps that are
earlier ends of crossing revisits of excursion greater than $\theta L$.
\end{corollary}

\begin{outline}[Outline of the formal proof of \cref{thm:X,cor:many}]
\emph{Thresholds.} From $d\ge d_0$ one gets $W\le\theta L$, so $1\le W\le H^*$, and $\theta L\ge2$, so $\theta L<H^*+1\le2H^*$.
\emph{Reduction.} Let $B$ separate $\Out(V_0)$ from $\CL_{H^*}$. If $\abs B\ge\abs{V_0}/2$ we are done, since
$\frac12\cE\le\frac1{24}$. Otherwise the set $V_0'$ of $q\in V_0$ with $\Out(q)\notin B$ has at least $\abs{V_0}/2\ge(\alpha/2)L$
elements, and $B$ separates it from $\CL_{H^*}$ within the depth $D_{\mathrm{E}}(M,C,\alpha/2)$. \emph{Bridge.} By the definition of
$\thetaM$, $H^*\le\theta_{\mathrm{E}}(k,D_{\mathrm{E}},C,\alpha/2)L$, and \cref{thm:bridge} at density $\alpha/2$ gives an acyclic
$U\subseteq V_0'$ in the ball relation with $\abs U\ge\rho_0\abs{V_0'}$. \emph{Coding.} \Cref{lem:acyclic-cost} with the genie of
cost $(\alpha\rho_0/8)L$ and depth $D_{\mathrm{E}}$ gives $\wval\abs B\ge\abs U-(\alpha\rho_0/8)L\ge(\rho_0/2)\abs{V_0'}$, hence
$\abs B\ge\cE\abs{V_0'}\ge\frac12\cE\abs{V_0}$.
\end{outline}

\begin{remark}[The earlier, row-scale form]\label{rem:2W}
With $H=W$ the threshold condition is automatic for large $d$, and the conclusion holds with $\CL_W$, the crossing steps of
revisits longer than two rows (\lean{conjectureE_obliviousLinTime_crossingLowerGT2W_menger_lb}); the same bound holds with the
larger set of both ends of revisits that cross or share a cell and a turning point with another revisit
(\lean{conjectureE_obliviousLinTime_mixGT2W_menger_lb}). Theorem~X is stronger: raising the threshold from $2W$ to $\theta L$
shrinks the target set, so rearrangements below a fixed fraction of $L$ can never substitute for macroscopic crossings.
\end{remark}

\begin{remark}[Numerical]\label{rem:numerics-crossings}
We computed exact minimum separating sets, by maximum flow, from large output sets to the earlier ends of crossing revisits of
excursion greater than $E$, for $E\in\{2W,L/64,L/16,L/4,L/2,L\}$, on several explicit oblivious schedules at
$d\in\{32,64,128\}$ and $b=1$. On the five linear-time cascade schedules of \cref{sec:cascades}, and on two superlinear transposers,
the ratio of the minimum separating set to $\abs{V_0}$ was exactly $1$ for every $E\le L/2$. At $E=L$ three of the cascades make no
such crossings at all, so the empty set separates, while heuristic bounds on their largest acyclic sets fell towards~$0$ as the
radius grew. So, numerically, the threshold in Theorem~D appears not to be a universal constant such as $\frac12$. These measurements are far below the
sizes at which the theorems apply and are no evidence for them; they suggest only that the true expansion constant may be of
order~$1$ at thresholds of order $L/C$.
\end{remark}

\section{Skeletons: where the crossings must be}\label{sec:skeletons}

Theorem~X says how much information must reach macroscopic crossings. This section says where they can be. At a scale that is
a fixed fraction of $L$, a linear-time head can turn only boundedly often, so its coarse motion is one of finitely many words;
crossings force two turns with overlapping legs; and the expansion of Theorem~X can always be found on one leg of one tape. A leg is
monotone only at the coarse scale: inside it the head may move back and forth over distances up to the scale.

\subsection{Crossings are zigzags}

\begin{lemma}[Crossing revisits are zigzags]\label{lem:zigzag}
\formal{MacroscopicSkeleton.lean}{crossing_isHop_opposite_polarity, exists_crossingZigzag_of_crossingHopsGT2H, card_tapeCrossingLowers_ge_of_macroOverlap}
Let $y$ be a unit-speed walk with horizon $T$ and $H\ge1$.
\begin{enumerate}[label=(\alph*)]
\item Two crossing revisits $(v_1,u_1)$, $(v_2,u_2)$ with $v_1<v_2<u_1<u_2$ visit different cells and lie on opposite sides of them:
one stays strictly above its cell and the other strictly below.
\item If two crossing revisits both have excursion greater than $2H$, there are a sign $\varepsilon\in\{\pm1\}$ and steps
$s_0<s_1<s_2<s_3<T$ such that $h=\varepsilon y$ satisfies
\[
  h(s_1)-h(s_0)\ge2H+1,\quad h(s_1)-h(s_2)\ge2H+1,\quad h(s_3)-h(s_2)\ge2H+1,\quad
  \min(h(s_1),h(s_3))-\max(h(s_0),h(s_2))\ge1 .
\]
\item Conversely, for every $\varepsilon$ and $s_0<s_1<s_2<s_3<T$, the walk has at least $2N$ steps that are earlier ends of
crossing revisits of excursion greater than $2H$, where $N=\max\bigl(0,\min(h(s_1),h(s_3))-\max(h(s_0),h(s_2))-2H\bigr)$, whatever the
walk does between these steps.
\end{enumerate}
\end{lemma}

\begin{outline}
(a) A walk with unit steps cannot pass a cell without visiting it (\lean{isWalk_ivt}); a revisit therefore stays on one side of its
cell, and the side of the second revisit is forced by the fact that the first cell is visited inside it but its own cell is not.
(b) Take the extreme points of the two excursions. (c) Every level strictly inside the overlap is crossed upward between $s_0$ and
$s_1$, downward between $s_1$ and $s_2$ and upward again between $s_2$ and $s_3$; the last visits before the turns give, for each
level, one upper and one lower revisit of excursion greater than $2H$, and these cross.
\end{outline}

\subsection{Skeletons}

\begin{definition}[Skeleton at scale $\Lambda$]\label{def:skeleton}
\formal{MacroscopicSkeleton.lean}{milestones, tapeSkeleton, macroTurns, traceSkeletonTuple}
Let $y$ be a walk with $y(0)=0$ and let $\Lambda\ge1$. Scan the steps $t=0,1,\dots,T-1$ keeping a grid point $g\in\Z$, initially
$0$. Whenever $\abs{y(t)-g\Lambda}\ge\Lambda$, record a \emph{milestone} $(t,\sigma)$ with $\sigma=+$ if $y(t)>g\Lambda$ and
$\sigma=-$ otherwise, and move $g$ one step in the direction~$\sigma$. The \emph{skeleton} of $y$ at scale $\Lambda$ is the word
$\sigma_1\sigma_2\cdots$ of recorded signs (\lean{tapeSkeleton}). A \emph{leg} is a maximal run of equal signs, and the number of
\emph{turns} is the number of sign changes (\lean{macroTurns}). The skeleton of a trace is the $k$-tuple of the skeletons of its
work tapes.
\end{definition}

\begin{lemma}[Confinement and finiteness]\label{lem:finite}
\formal{MacroscopicSkeleton.lean}{epochGrid_confines_and_exact, tapeSkeleton_length_mul_le, card_boundedSkeletonTuple, traceSkeletonTuple_mem_boundedSkeletonTuple}
\begin{enumerate}[label=(\alph*)]
\item At every step $\abs{y(t)-g(t)\Lambda}<\Lambda$, where $g(t)$ is the current grid point, and $y(t)=g(t)\Lambda$ exactly at
each milestone. A skeleton of length $m$ has $m\Lambda\le T$.
\item Let $\tau$ be any trace with $k$ work tapes and $T\le C(L+1)$, with $d,b\ge1$, $C\ge1$, $\theta>0$ and $d\ge\lceil8/\theta\rceil$.
At scale $\Lambda=\lfloor\theta L/4\rfloor$ we have $\Lambda\ge1$, and every work tape's skeleton has at most
$\Mmax=\lfloor16C/\theta\rfloor$ letters. Hence the skeleton of $\tau$ belongs to the set of $k$-tuples of words of length at most
$\Mmax$ over $\{+,-\}$, which has exactly $(2^{\Mmax+1}-1)^k$ elements and does not depend on~$d$.
\end{enumerate}
\end{lemma}

\begin{outline}
(a) is an induction over the scan. (b) Each milestone after the first needs the head to move $\Lambda$ cells, so $m\Lambda\le T\le
2CL$, while $\Lambda\ge\theta L/8$ once $\theta L\ge8$; hence $m\le16C/\theta$. The count of words of length at most $\Mmax$ is
$\sum_{i\le\Mmax}2^i$.
\end{outline}

\begin{definition}[Crossing-capable skeletons]\label{def:capable}
\formal{MacroscopicSkeleton.lean}{IsScaleFreeCrossingCapable, CrossingCapableSkeletonTuple}
Let $s$ be a skeleton and let $\gamma(i)$ be its grid coordinate after $i$ letters. The skeleton is \emph{crossing-capable} if for
some sign $\varepsilon$ and indices $i_0<i_1<i_2<i_3\le\abs s$ the function $h=\varepsilon\gamma$ satisfies
\[
  h(i_1)-h(i_0)\ge3,\quad h(i_1)-h(i_2)\ge3,\quad h(i_3)-h(i_2)\ge3,\quad \min(h(i_1),h(i_3))-\max(h(i_0),h(i_2))\ge-1
\]
(\lean{IsScaleFreeCrossingCapable}): up three grid steps, down three, up three, with overlapping ranges. Such a skeleton has at
least two turns.
\end{definition}

\begin{proposition}[Which skeletons carry crossings]\label{prop:capable}
\formal{MacroscopicSkeleton.lean}{isCrossingCapableSkeleton_of_tapeCrossingLowers_nonempty, isScaleFreeCrossingCapable_of_two_mul_le, two_le_macroTurns_of_isScaleFreeCrossingCapable, tapeCrossingLowers_nonempty_of_isDetourProofSkeleton}
Let $y$ be a walk with $y(0)=0$, $\Lambda\ge1$ and $H\ge2\Lambda$. If the tape of $y$ has an earlier end of a crossing pair of revisits
of excursion greater than $2H$, then its skeleton at scale $\Lambda$ is crossing-capable; in particular it has at least nine
letters and at least two turns. Conversely, if for some sign and indices $i_0<i_1<i_2<i_3$ the skeleton's grid satisfies
$(\min(h(i_1),h(i_3))-\max(h(i_0),h(i_2)))\Lambda\ge2H+1$, then the tape has such crossing steps.
\end{proposition}

\begin{outline}
The first part combines \cref{lem:zigzag}(b) with the confinement of \cref{lem:finite}(a): the four extreme steps lie within
$\Lambda$ of grid points, which turns the inequalities with margin $2H+1\ge4\Lambda+1$ into the grid inequalities above
(\lean{isCrossingCapableSkeleton_of_tapeCrossingLowers_nonempty}, which proves a finer statement at every scale). The converse is
\cref{lem:zigzag}(c) applied at the milestone steps, where the walk is exactly on the grid.
\end{outline}

Between the two conditions there is a narrow band of skeletons that may or may not carry crossings, depending on the motion inside
the legs; this band is non-empty, so no condition on the skeleton alone decides whether crossings occur.

\begin{lemma}[No macroscopic revisit within a leg]\label{lem:legs}
\formal{MacroscopicSkeleton.lean}{legOfEpoch_lt_of_isHop_gt2H}
Let $y$ be a walk with $y(0)=0$, $\Lambda\ge1$ and $H\ge2\Lambda$. Every revisit of excursion greater than $2H$ starts in an earlier
leg of the skeleton than the one in which it ends.
\end{lemma}

\begin{outline}
Within one leg the grid point moves monotonically. A revisit returns to its cell, so if it started and ended in the same leg it
would be confined to at most two adjacent epochs, with excursion less than $3\Lambda\le2H$.
\end{outline}

\subsection{Theorem S}

For an oblivious linear-time transposer $M$ with time constant $C$ and $\alpha>0$, put $\theta=\thetaM(M,C,\alpha)$,
$\Lambda^*=\lfloor\theta L/4\rfloor$, $H^*=\lfloor\theta L\rfloor$ and $\Mmax=\lfloor16C/\theta\rfloor$, and for a work tape $j$ and
$r\ge0$ let $\CL^{j,r}$ be the set of steps of $\CL_{H^*}$ that are earlier ends of revisits on tape $j$ and lie in the $r$-th leg of
tape $j$'s skeleton at scale $\Lambda^*$ (\lean{legCrossingLowers}).

\begin{theorem}[Theorem~S, skeletons and a single leg]\label{thm:S}
\formal{MacroscopicSkeleton.lean}{obliviousLinTime_reduces_to_crossingCapableSkeleton_and_singleLeg, obliviousLinTime_infinite_pinned_skeleton_and_leg}
Let $M$ be a machine with $k$ work tapes, $C\ge1$ and $0<\alpha\le1$, and let $d\ge\max(d_0,\lceil8/\theta\rceil)$ with $d_0$ as in
\cref{thm:X}. Suppose $\mathrm{ObliviousLinTime}(M,d,C,T)$.
\begin{enumerate}[label=(\alph*)]
\item $1\le\Lambda^*$ and $2\Lambda^*\le H^*$, and the skeleton of the common trace at scale $\Lambda^*$ belongs to the explicit
finite family of $k$-tuples of words of length at most $\Mmax$ in which some tape's word is crossing-capable.
\item For every set $V_0$ of at least $\alpha L$ input positions there are a work tape $j$, whose skeleton is crossing-capable
with at least nine letters and two turns, and a leg $r\le\Mmax$ such that every set separating $\Out(V_0)$ from $\CL^{j,r}$ has at
least
\[
  \frac{\cE(M,C,\alpha/2)}{2k(\Mmax+1)}\,\abs{V_0}
\]
elements.
\item If $\mathrm{ObliviousLinTime}(M,d,C,T_d)$ holds for all $d$ in an infinite set $S$, then there are one tuple of words, one tape
and one leg such that, for infinitely many $d\in S$, the skeleton is that tuple and (b) holds with that tape and leg for $V_0$ the
set of all input positions (with $\alpha=1$).
\end{enumerate}
\end{theorem}

\begin{outline}
(a) is \cref{lem:finite}(b) together with \cref{prop:capable}, since $\CL_{H^*}$ is non-empty by \cref{cor:many}. (b) By
\cref{lem:legs} each step of $\CL_{H^*}$ lies in a slot $(j,r)$ with $r\le\Mmax$, and there are at most $k(\Mmax+1)$ slots. If every
slot had a separating set smaller than the stated bound, their union would separate $\Out(V_0)$ from all of $\CL_{H^*}$
(\lean{isMengerVertexCut_biUnion}) with fewer than $\frac12\cE\abs{V_0}$ elements, contradicting \cref{thm:X}. (c) The skeleton,
tape and leg take finitely many values, independent of $d$, so one of them recurs infinitely often
(\lean{exists_infinite_fiber_of_mapsTo_finset}).
\end{outline}

The tape and the leg in (b) may depend on $V_0$ and on~$d$; (c) removes the dependence on $d$ along a subsequence. The finite
family in (a) is astronomically large, because $\Mmax=\lfloor16C/\theta\rfloor$ inherits the size of~$\theta$; what matters is that
it does not grow with~$d$. The Lean proves only that the crossing-capable tuples form a proper subfamily of all tuples (the idle
tuple is excluded, \lean{card_crossingCapableSkeletonTuple_lt}), not how much smaller it is.

\begin{remark}[Numerical: what the skeleton does not see]\label{rem:skeleton-limits}
Inside a leg the head may move back and forth over distances up to $\Lambda^*$, which is a fixed fraction of $L$. The skeleton does
not see this motion, and the explicit cascades of \cref{sec:cascades} hide all but boundedly many of their refinement stages inside
it. Our measurements on those schedules, and on a one-pass shear that \cref{thm:J} excludes, found skeletons independent of $d$ with a
single leg carrying linear expansion ($d\le128$). So, as far as these measurements go, the skeleton condition does not separate these
schedules from a transposer.
\end{remark}

\section{Rows to columns}\label{sec:demand}

Theorems~H, X and~S use transposition only through the certificate method: a transposer cannot have a large acyclic set after
a cheap cut. This section uses the transposition demand directly. Information from a set of input rows must reach the output rows
into which transposition sends it, and every cut between the two must carry it.

For sets $A_0,A_1\subseteq\{0,\dots,d-1\}$ of rows let $\Lin(A_0)$ be the set of steps at which the input head reads a position in
an input row of $A_0$ (\lean{inSlabReadSteps}), and let $\Qstar(A_0,A_1)$ be the set of input positions whose input row lies in
$A_0$ and whose output row lies in $A_1$ (\lean{bichromaticRect}): the bits that rows $A_0$ send into output rows $A_1$. It has
exactly $\abs{A_0}\abs{A_1}b$ elements (\lean{card_bichromaticRect}).

\begin{lemma}[Lemma~R, rows to columns]\label{lem:R}
\formal{TranspositionCut.lean}{isAcyclicIn_coneSystem_uncutRootedBichromaticRect, frozenTrace_bichromatic_tposOf_cut_ge, oblivious_bichromatic_tposOf_cut_ge}
Let $M$ transpose at size $d\ge1$ within time $c$, and let $\tau$ be an $\eta$-heavy trace of $M$; for an oblivious $M$ take its common
trace and $\eta=0$. For all sets $A_0,A_1$ of rows and every set $\Gamma$ of steps separating $\Out(\Qstar(A_0,A_1))$ from $\Lin(A_0)$,
\[
  \wval\abs\Gamma+\abs{\Gamma\cap\Out(\Qstar(A_0,A_1))}\ \ge\ \abs{A_0}\abs{A_1}b-\eta L,
  \qquad\text{hence}\qquad \abs\Gamma\ \ge\ \frac{\abs{A_0}\abs{A_1}b-\eta L}{\wval+1}.
\]
\end{lemma}

\begin{outline}
Let $Q'$ be the set of $q\in\Qstar(A_0,A_1)$ whose root is not in $\Gamma$; since $\Out$ is injective,
$\abs{Q'}\ge\abs{A_0}\abs{A_1}b-\abs{\Gamma\cap\Out(\Qstar)}$. For $q\in Q'$, a dependency path from $\Out(q)$ that avoids $\Gamma$ after
its root cannot end at a read of rows $A_0$: the root is not in $\Gamma$, so the path would avoid $\Gamma$ altogether. Hence the cone of
every $q\in Q'$ avoiding $\Gamma$ contains no position of rows $A_0$, while $Q'$ lies in rows $A_0$. No element of $Q'$ sees another, so
$Q'$ is acyclic, and \cref{lem:frozen} gives $\abs{Q'}-\wval\abs\Gamma\le\eta L$.
\end{outline}

With $A_0=A_1$ the set of all rows, Lemma~R says that every cut between the input and the output carries the whole input:
$\abs\Gamma\ge(1-\eta)L/(\wval+1)$. With smaller row sets it says the same for every block of the transposition demand separately. The
filtering of roots matters: $\Gamma=\Out(\Qstar)$ is a separating set, and the cone relation avoiding it need not be acyclic on all of
$\Qstar$, because roots in $\Gamma$ may still see their own inputs.

The lemma becomes a statement about a single set of steps through a colouring.

\begin{theorem}[Colouring criterion]\label{thm:colouring}
\formal{TranspositionCut.lean}{isMengerVertexCut_of_bichromaticColoring, frozenTrace_bichromaticColoring_wave_cut_ge}
Let $M$, $d$, $c$, $\eta$ and $\tau$ be as in \cref{lem:R}, with halting time $T$, and let $A_0,A_1$ be sets of rows. Let $Y$ be a set of
steps and $\chi$ a colouring of the steps by $\{0,1\}$ such that
\begin{enumerate}[label=(\roman*)]
\item $\chi=0$ on $\Lin(A_0)\setminus Y$ and $\chi=1$ on $\Out(\Qstar(A_0,A_1))\setminus Y$;
\item every revisit $(u,t)$ with $u,t\notin Y$ has $\chi(t)\le\chi(u)$.
\end{enumerate}
Let $\partial\chi$ be the set of steps $t\in[1,T)$ with $t,t-1\notin Y$, $\chi(t)=1$ and $\chi(t-1)=0$. Then every set $B$ of steps separating
the roots of all input positions from $Y$ satisfies
\[
  \abs B\ \ge\ \frac{\abs{A_0}\abs{A_1}b-\eta L}{\wval+1}-\abs{\partial\chi}.
\]
\end{theorem}

\begin{outline}
$B\cup\partial\chi$ separates $\Out(\Qstar)$ from $\Lin(A_0)$. Take a dependency path from a root to a read of rows $A_0$ that avoids
$B\cup\partial\chi$. It avoids $Y$, since its part up to a vertex of $Y$ would be a path from a root to $Y$ avoiding $B$. By (i) the path
starts with colour $1$ and ends with colour $0$, so some edge drops from colour $1$ to colour $0$ as the path moves to an earlier step. A
state edge $t\to t-1$ with $\chi(t)=1$ and $\chi(t-1)=0$ would put $t$ in $\partial\chi$, and a revisit edge $t\to u$ with $\chi(t)=1$ and
$\chi(u)=0$ is excluded by (ii), a contradiction. \Cref{lem:R} bounds $\abs{B\cup\partial\chi}$.
\end{outline}

Taking $B=Y$ shows $\abs Y\ge(\abs{A_0}\abs{A_1}b-\eta L)/(\wval+1)-\abs{\partial\chi}$: a set of steps across which the colour can only
jump at few places must itself be large. Several such sets that are pairwise disjoint give a time bound. For $d\ge16$ let $K=\lfloor
\log_2d\rfloor$ and $d_1=2^K$, and for $0\le s\le K-2$ let $A_s^0$ and $A_s^1$ be the rows $i<d_1$ whose bit $K-1-s$ is $0$, respectively $1$.

\begin{corollary}[Scale-separated traces]\label{cor:scale-separated}
\notformalized
Let $M$, $d$, $c$, $\eta$ and $\tau$ be as in \cref{lem:R}, with $d\ge16$, $\eta\le\frac1{32}$, $K_{\mathrm{flip}}\ge0$ and
$d\ge64(\wval+1)K_{\mathrm{flip}}$. Suppose that for each $s$ in a set $S\subseteq\{0,\dots,K-2\}$ there are a set $Y_s$ of steps and a
colouring $\chi_s$ satisfying (i) and (ii) of \cref{thm:colouring} for $(A_s^0,A_s^1)$ with $\abs{\partial\chi_s}\le K_{\mathrm{flip}}d$, and
that the sets $Y_s$ are pairwise disjoint. Then $T\ge\abs S\cdot L/(64(\wval+1))$.
\end{corollary}

\begin{informalproof}
Since $d_1>d/2$, $\abs{A_s^0}\abs{A_s^1}b=(d_1/2)^2b>L/16$. Since $d\ge64(\wval+1)K_{\mathrm{flip}}$ and $b\ge1$,
$K_{\mathrm{flip}}d\le d^2/(64(\wval+1))\le L/(64(\wval+1))$. \Cref{thm:colouring} with $B=Y_s$ gives
\[
  \abs{Y_s}\ \ge\ \frac{(1/16-\eta)L}{\wval+1}-K_{\mathrm{flip}}d\ \ge\ \frac{L}{32(\wval+1)}-\frac{L}{64(\wval+1)}=\frac{L}{64(\wval+1)} .
\]
The $Y_s$ are pairwise disjoint subsets of $\Range T$, so $T\ge\sum_{s\in S}\abs{Y_s}$.
\end{informalproof}

With $S$ all $K-1$ scales the bound is $T\ge(\lfloor\log_2d\rfloor-1)L/(64(\wval+1))$, so no linear-time transposer has a trace that is
scale-separated across all scales. The corollary should be read with \cref{prop:onecut} in mind: one set of steps, together with the
colouring that changes once, at the step after the last input read, satisfies (i) and (ii) for every pair of row sets at once, so the content
of the hypothesis is the disjointness of the sets $Y_s$, which says roughly that live storage is rewritten once per scale. The corollary is
therefore a time bound for schedules that refine the row index one bit at a time in separate passes, whatever they compute; it does not
show that a linear-time transposer must do so.

\section{Adaptive machines}\label{sec:adaptive}

An adaptive machine has many traces and no common dependency DAG. Two routes lead from the oblivious results to adaptive machines.
The first uses a single heavy trace of the machine itself; the second passes through the bit simulator of the fifteenth paper, whose
traces are few.

\begin{theorem}[Theorem~E, heavy traces and simulators]\label{thm:E}
\formal{TightConstants.lean}{frozenTrace_crossingLowerGT2H_menger_lb_tight, frozenTrace_macroscopicCrossingLowers_menger_lb_tight, frozenTrace_reduces_to_crossingCapableSkeleton_and_singleLeg_tight, simulator_heavyTrace_expansion_and_skeleton_tight, adaptive_simulator_heavyTrace_macroscopicExpansion_tight}
\begin{enumerate}[label=(\alph*)]
\item \emph{Heavy traces.} Let $M$ be any machine with $k$ work tapes, $C\ge1$ and $\alpha>0$, put $\theta=\thetaM(M,C,\alpha)$ and
$\rho_0=\rho_0(k,C,\tfrac{19}{20}\alpha)$, and let $d\ge\max\bigl(1,d_{\mathrm{br}}(k,C,\tfrac{19}{20}\alpha)\bigr)$. Suppose $M$ transposes at size
$d$ within time $C$ (as it does if it transposes within a smaller time), and let $\tau$ be an $\eta$-heavy trace of $M$. Then for every
$1\le H\le\theta L$ and every set $V_0$ of at least $\alpha L$ input positions, every set separating $\Out(V_0)$ from $\CL_H$ in the
dependency DAG of $\tau$ has at least
\[
  \cE(M,C,\alpha,\eta)\,\abs{V_0},\qquad \cE(M,C,\alpha,\eta)=\frac{\rho_0/2-\eta/\alpha}{\wval},
\]
elements; at $\eta=0$ this is $\cE(M,C,\alpha)$. If moreover $\eta\le(\alpha/8)\rho_0$ and $d\ge d_0$ with $d_0$ as in \cref{thm:X}, the
conclusion of \cref{thm:X} at $H^*$ holds in the dependency DAG of $\tau$ with the expansion constant
$\cE'(M,C,\alpha)=\frac38\rho_0/\wval\le\cE(M,C,\alpha,\eta)$ in place of $\cE(M,C,\alpha)$; if in addition $\alpha\le1$ and
$d\ge\lceil(8/\theta)^{1/2}\rceil$, so do the conclusions (a) and (b) of \cref{thm:S}, with $\cE'$ in place of $\cE$ in (b).
\item \emph{Simulators.} Let $M$ be any machine with $k$ work tapes, $c>0$ and $0<\alpha\le1$, and put $\eta^*=(\alpha/16)\rho_0(k,11c+1,\tfrac{19}{20}\alpha)$.
There are $s\ge1$ and $d_{\mathrm{sim}}$ such that the simulator $\Ns$ of~\cite[Theorem~S]{HighRate2026}, which has the same work tapes and
work alphabet as $M$ and computes the same partial function, satisfies the following at every $d\ge d_{\mathrm{sim}}$ at which $M$
transposes within time $c$: $\Ns$ transposes within time $11c+1$; it has at most $2^{\eta^*L}$ distinct traces on the inputs of length
$L$; its $2\eta^*$-heavy traces are followed together by more than $(1-2^{-\eta^*L})2^L$ inputs; and on each of its $2\eta^*$-heavy traces,
with $\theta$ and $H^*$ computed for $\Ns$ and the time constant $11c+1$, $W\le H^*$, $\theta L\le2H^*$, the separation bound of \cref{thm:X} at
$H^*$, and the conclusions of \cref{thm:S}(a),(b) all hold with the expansion constant $\cE'$ of $\Ns$.
\end{enumerate}
\end{theorem}

\begin{outline}
(a) \Cref{lem:frozen} gives $\abs U-\wval\abs\Gamma\le\eta L$ for every set $\Gamma$ of steps and every set $U$ acyclic in the cone
relation of $\tau$ avoiding $\Gamma$. The proof of \cref{thm:X} uses the transposition hypothesis only through such a bound, with $0$ in
place of $\eta L$, and gives $\wval\abs B\ge\frac12\rho_0\abs{V_0}-\eta L\ge(\frac12\rho_0-\eta/\alpha)\abs{V_0}$, since
$\alpha L\le\abs{V_0}$ (\lean{linearTimeTrace_crossingLowerGT2H_menger_lb_of_coneDeficit_tight}). For $\eta\le(\alpha/8)\rho_0$ this is at
least $\frac38\rho_0\abs{V_0}$. (b) The constant $\eta^*$ depends only on $k$, $c$
and $\alpha$. The fifteenth paper's trace transfer at $\eta^*$ (\lean{TapeStorageRates.aligned_low_trace_dimensions}) gives $s$ and the
simulator, with at most $2^{\eta^*L}$ traces; the light traces are followed by fewer than $2^{\eta^*L}2^{(1-2\eta^*)L}$ inputs. Since $\Ns$ has
$k$ work tapes and runs in time $11c+1$, and $2\eta^*=(\alpha/8)\rho_0(k,11c+1,\tfrac{19}{20}\alpha)$, part (a) applies to $\Ns$ on each heavy
trace.
The quantifiers are ordered $\eta^*$, then $s$ and $\Ns$, then the constants of $\Ns$ and $d_{\mathrm{sim}}$.
\end{outline}

Part (b) is the only route by which the macroscopic condition reaches arbitrary linear-time transposers so far, and it reaches them
through a different machine's dependency DAG. The simulator rewrites windows of $3s$ cells in each round, and a separating set in its
DAG need not correspond to one in the DAG of $M$ (\cref{prob:pullback}). Its value width grows with $s$, so its constants are much smaller
than those of $M$. Statements in this paper about the dependency DAG of a general machine therefore always assume a heavy trace or name the
simulator.

\begin{remark}[The heaviness hypothesis cannot simply be dropped]\label{rem:rare}
Proof sketch only: given any machine $M_0$, add three states so that the new machine first scans the input for a $1$; on the all-zero
input it writes the all-zero output with its work heads standing still, and otherwise it returns its input head and runs $M_0$. The
new machine transposes whenever $M_0$ does, within a constant more time, and its all-zero trace makes no revisits at all. Hence the
version of \cref{thm:E}(a) without a lower bound on the trace probability would exclude this machine, and so would imply $\neg$(T).
It is not refuted: it is false if (T) holds.
\end{remark}

\begin{remark}[Wave counts on heavy traces]\label{rem:waves}
Because every expanding set of steps is as large as its expansion, $m$ pairwise disjoint sets of steps into which a set of $\alpha L$ outputs
expands by a factor $c$ need time at least $mc\alpha L$ (\lean{disjointWave_mengerExpansion_counting_le}); for oblivious machines in time
$C(L+1)$ there are at most $2C/(c\alpha)$ of them (\lean{obliviousLinTime_expandingWaves_le}). A hypothesis asserting unboundedly many such
sets on every heavy trace of every linear-time transposer is therefore equivalent to there being no linear-time transposer at
infinitely many sizes, and gives no route to $\neg$(T). For oblivious machines, a hypothesis asserting unboundedly many such sets on every
transposer, of any running time, would exclude oblivious linear-time transposers; the converse is not known.
\end{remark}

The fibre argument needs correctness only on the inputs it counts, so \cref{thm:X} and the separation bound of \cref{thm:E}(a)
extend to machines that transpose correctly at least a $2^{-\eta L}$ fraction of their inputs. The constant $\cE(M,C,\alpha,\eta)$
of \cref{thm:E}(a) equals $\cE(M,C,\alpha)-\eta/(\alpha\wval)$, and it is positive exactly for
$\eta<\frac12\alpha\rho_0(k,C,\tfrac{19}{20}\alpha)$.

\begin{corollary}[Approximate transposers]\label{cor:approx}
\formal{TightConstants.lean}{correctSubset_crossingLowerGT2H_menger_lb_tight, obliviousOn_correctSubset_macroscopicCrossingLowers_menger_lb_tight}
Let $M$ be any machine with $k$ work tapes, $C\ge1$, $\alpha>0$ and $\eta\ge0$.
\begin{enumerate}[label=(\alph*)]
\item Let $d\ge\max\bigl(1,d_{\mathrm{br}}(k,C,\tfrac{19}{20}\alpha)\bigr)$, and let $\tau$ be a trace of $M$ followed by at least
$2^{(1-\eta)L}$ inputs that $M$ transposes within time $C$. For every $1\le H\le\thetaM(M,C,\alpha)L$ and every set $V_0$ of at least
$\alpha L$ input positions, every set of steps separating $\Out(V_0)$ from $\CL_H$ in the dependency DAG of $\tau$ has at least
$\cE(M,C,\alpha,\eta)\abs{V_0}$ elements.
\item Let $d\ge d_0$ with $d_0$ as in \cref{thm:X}, and let $M$ be oblivious at length $L$ with halting time $T$ and transpose at least
$2^{(1-\eta)L}$ inputs within time $C$. Then the conclusion of \cref{thm:X} at $H^*$ holds in the common trace, with $\cE(M,C,\alpha,\eta)$
in place of $\cE(M,C,\alpha)$; for other $H$, (a) applies, since every input follows the common trace.
\end{enumerate}
\end{corollary}

\begin{outline}
(a) The proof of \cref{thm:E}(a) goes through with \cref{lem:frozen} applied to the correctly transposed inputs that follow $\tau$; the
deficit $\eta L\le(\eta/\alpha)\abs{V_0}$ lowers the constant by $\eta/(\alpha\wval)$. (b) All inputs follow the common trace, which halts
at $T$, and any correctly transposed input gives $T\le C(L+1)$; apply (a) at $H^*$. With $\eta=0$, (b) is \cref{thm:X}.
\end{outline}

Read the other way, a certificate on each of few traces bounds the fraction of inputs that a machine transposes correctly.

\begin{corollary}[Deficient traces carry few correct inputs]\label{cor:frozen-count}
\formal{SimulatorBridge.lean}{card_correctInputs_le_of_coneDeficit}
Let $M$ be any machine, $d\ge1$, $c>0$, and $\mathcal T$ a finite set of traces of $M$, each with a \emph{deficit certificate}: a set
$\Gamma_\tau$ of steps and a set $U_\tau$ acyclic in the cone relation of $\tau$ avoiding $\Gamma_\tau$ with
$\abs{U_\tau}-\wval\abs{\Gamma_\tau}\ge\delta L$. Then at most $\abs{\mathcal T}\,2^{(1-\delta)L}$ inputs follow a trace in $\mathcal T$
and are transposed by $M$ within time $c$. If $\abs{\mathcal T}\le2^{\eta L}$, these inputs are at most a $2^{-(\delta-\eta)L}$ fraction of
all inputs.
\end{corollary}

\begin{outline}
For each $\tau\in\mathcal T$ followed by correctly transposed inputs, \cref{lem:frozen} applied to the set $S_\tau$ of these inputs gives
$\delta L\le L-\log_2\abs{S_\tau}$. Sum $\abs{S_\tau}\le2^{(1-\delta)L}$ over $\mathcal T$.
\end{outline}

\section{Class exclusions}\label{sec:classes}

The results of this section exclude classes of machines defined by their head motion. They use the certificate method directly and
reach machines that Theorems~X and~S do not: schedules with macroscopic crossings that are nevertheless excluded, and adaptive machines
with no heavy trace.

\subsection{Boundedly many macroscopic phases}

A walk is \emph{weakly monotone} on an interval of steps if it never moves left there, or never moves right (pauses allowed). It is
\emph{$A$-coarsely monotone} if it never moves more than $A$ cells against a fixed direction: $x(t')\ge x(t)-A$ for all $t\le t'$ in the
interval, or $x(t')\le x(t)+A$ for all such $t\le t'$ (\lean{CoarseMonoOnInterval}).

\begin{definition}[Bounded macro-phases]\label{def:BM}
\formal{MacroBankCrux.lean}{WeakMonoOnInterval, BoundedMacroBank}
Let $M$ be a machine, $C\ge1$, $M_{\mathrm{ph}}\ge1$, $B$ a set of work heads (the \emph{bank}) and $S:\N\to\N$ with $S(d)/L\to0$. $M$ has
\emph{bounded macro-phases} with these parameters (\lean{BoundedMacroBank}) if for all large $d$, whenever $M$ is oblivious at length $L$
with halting time $T\le C(L+1)$, the interval $\Range T$ splits into $M_{\mathrm{ph}}$ consecutive intervals on each of which the input
head and every bank head are weakly monotone, and every other work head visits at most $S(d)$ cells before~$T$.
\end{definition}

The heads outside the bank are unrestricted inside their $o(L)$ workspaces: any number of reversals, any nesting, any coding.

\begin{theorem}[Theorem~B, bounded macro-phases]\label{thm:BM}
\formal{SuffixMassCrux.lean}{boundedMacroBank_has_crux, boundedMacroBank_excludes, suffixMacroBank_has_crux, suffixMacroBank_excludes}
If $M$ has bounded macro-phases with bank $B$ and $M_{\mathrm{ph}}$ phases, then $M$ has crux rate $1/(4(\abs B+1)^{M_{\mathrm{ph}}-1})$ at every
depth, and for all large $d$ there is no $T$ with $\mathrm{ObliviousLinTime}(M,d,C,T)$. The same holds, with crux rate
$\delta_0/(4(\abs B+1)^{M_{\mathrm{ph}}-1})$, if $\delta_0>0$ and the conditions of \cref{def:BM} hold only on a final interval of steps that
contains every append, with at least $\delta_0d$ input rows not read before that interval; the motion before it is unrestricted.
\end{theorem}

\begin{outline}
We prove the suffix form; the whole-run form is its case $\delta_0=1$, with the whole run as the final interval
(\lean{boundedSuffixMacroBank_of_boundedMacroBank}). Dependency paths decrease time, so a path from an output to a read of an untouched
row stays inside the final interval (\lean{getLastD_le_of_mem_depPath}). Let $K_0=(\abs B+1)^{M_{\mathrm{ph}}-1}$ and cut workspace revisits
in the final interval longer than $\lceil\sqrt{S(T+1)}\rceil$; there are $o(L)$ of them. Inside a phase a bank revisit goes back by one step;
between phases, bank revisits are strictly monotone partial injections of fixed sign. So after deleting $o_D(L)$ exceptional outputs,
every short dependency path from an output has a \emph{signature}, the decreasing sequence of phases it enters through bank heads, and at
most $K_0$ signatures start in a given phase. Thresholds on input addresses pull back along signatures to thresholds in time. Split the
untouched input rows into $2K_0$ buckets: each unsplit output row forbids at most $K_0$ buckets, so some bucket is missed by half of the
output rows, and the rectangle of that bucket and those output rows, minus the exceptions, is independent in the cone relation; it has
$(\delta_0/(4K_0)-o(1))L$ elements (\lean{suffixMacroBank_independent_exists}).
\end{outline}

For adaptive machines the phases and the workspaces may depend on the input, and the walks may jitter.

\begin{theorem}[Theorem~J, jittered phases, any adaptivity]\label{thm:J}
\formal{AdaptiveSuffixJitterBM.lean}{native_machine_preserves_suffixJitterBM, nativeGapBound_of_adaptiveJitterBM, adaptiveJitterBM_excludes_transposesAt, adaptiveJitterBM_not_linearTimeTransposer, adaptiveSuffixJitterBM_excludes_transposesAt, adaptiveSuffixJitterBM_not_linearTimeTransposer}
Let $M$ be any machine, $c>0$, $M_{\mathrm{ph}}\ge1$, $B$ a set of work heads, and $A(d),S(d)=o(L)$. Suppose that for all large $d$ every run of $M$ on an input of length $L$ that halts
within $c(L+1)$ steps splits into $M_{\mathrm{ph}}$ intervals on each of which the input head and the bank heads are $A(d)$-coarsely
monotone, while every other work head visits at most $S(d)$ cells; the intervals and cells may depend on the input. Then $M$ does not
transpose within time $c$ at all large $d$; if such $M_{\mathrm{ph}}$, $B$, $A$ and $S$ exist for every $c\ge1$, then $M$ is not a linear-time transposer. The same holds if $\delta_0>0$ and the conditions hold only
on a final interval containing every append with at least $\delta_0d$ rows not entered by the input head before it.
\end{theorem}

\begin{outline}
The bit simulator of the fifteenth paper preserves the shape, also in its suffix form, with the jitter and the workspace enlarged by
$O(s)$ (\lean{native_machine_preserves_suffixJitterBM}). On each frozen trace of the simulator, the proof of Theorem~B goes through with
jitter: the revisits that straddle a pulled-back cut on a bank head occupy at most $2A+1$ cells (\lean{hop_upper_injOn}), and input reads
are pulled back up to a band of width~$A$. This gives every storage code on the frozen cone relations of every simulator rate at most
$1-1/(4K_0)+o(1)$ at the sizes at which $M$ transposes within time~$c$ (\lean{nativeGapBound_of_adaptiveJitterBM}), which is the hypothesis
of \cref{thm:G} below; so Theorem~G excludes~$M$.
\end{outline}

At a fixed number of phases these classes, in their whole-run forms, also yield to pass counting at their own time bounds, by
\cref{thm:coarsepass}, which extends the grouped pass bound of the first paper~\cite[Theorems~5.8 and~5.12]{FirstPaper2026} to coarse
passes with small workspaces. Moreover, the same theorem excludes these shapes, imposed on whole runs of any length, at any running time
for oblivious machines and at any polynomial running time for adaptive ones, even with jitter and workspaces up to a small fixed fraction
of $L$ (\cref{cor:boundedpasses}). What Theorems~B and~J add is an explicit native coding
deficit, and the suffix forms, which pass counting does not reach: an unrestricted prefix may read $(1-\delta_0)d$ rows, and conditioning on the configuration at the start of the suffix
leaves classes of inputs of co-entropy up to $(1-\delta_0)L$, for which the published robust pass bounds~\cite{FirstPaper2026,PassesMemo2026}
say nothing.

\subsection{A native gap suffices}

\paragraph{Roots.} The hypotheses of \cref{thm:G} and \cref{cor:G-uniform,cor:G-low} below are those of the fourteenth paper's native
bound~\cite[Corollary~L]{BoundedStack2026} with a general deficit (\lean{NativeGapBound}). They concern the cone relation of a run
\emph{with arbitrary roots}: for each output position, the root may be any step before the halting step at which the output head is on
that position's cell, and the rate bound is required for every such choice of roots. The append steps are one such choice; whether the
results hold when the bound is required only for them is not proved.

\begin{theorem}[Theorem~G, native gap criterion]\label{thm:G}
\formal{NativeGapCriterion.lean}{nativeGapBound_half_iff, native_gap_excludes_transposesAt, native_gap_excludes_linear}
Let $M$ be a machine and $c>0$. Suppose that there is $\rho>0$ such that, for every $s\ge1$ and every depth $D$, every binary storage
code on the cone relation with arbitrary roots, after a genie of full depth $D$, of every run of the simulator $\Ns$ of $M$~\cite[Theorem~S]{HighRate2026}
that halts within $(11c+1)(L+1)$ steps on an input of length $L$ (entries of $b(d)$ bits) has rate at most $1-\rho+o(1)$, where $d$ ranges
over the sizes at which $M$ transposes within time~$c$. Then $M$ transposes within time $c$ at only finitely many $d$. If such a $\rho$
exists for every machine $M$ and every $c>0$, then (T) fails.
\end{theorem}

\begin{outline}
Choose $s$ so that, at all large $d$ at which $M$ transposes within time $c$, $\Ns$ transposes within time $11c+1$ with at most
$2^{(\rho/2)L}$ traces~\cite[Theorem~S]{HighRate2026}. Some trace of $\Ns$ carries $2^{(1-\rho/2)L}$ inputs; after the genie, of cost
$\frac14\rho L$, its fibre splits into storage codes on its cone relation (\lean{MergeDepth.blockMachine_isStorageCode_fibre}), one of
which has rate at least $1-\frac34\rho-o(1)$, contradicting the hypothesis. If (T) held, some machine $M$ would transpose in time
$c(L+1)$ at infinitely many sizes. At $\rho=\frac12$ the hypothesis is the native half bound of the fourteenth
paper~\cite[Corollary~L]{BoundedStack2026} on the simulators of~$M$.
\end{outline}

\begin{corollary}[Theorem~G, uniform form]\label{cor:G-uniform}
\formal{NativeGapCriterion.lean}{native_gap_exists_exclude_linear, native_gap_tests_exclude_linear}
Suppose that for every $k$ and $C\ge1$ there is $\rho(k,C)>0$ such that, for every machine with $k$ work tapes and every depth $D$, every
binary storage code on the cone relation with arbitrary roots, after a genie of full depth $D$, of every run that halts within $C(L+1)$ steps on an input of
length $L$ (entries of $b(d)$ bits) has rate at most $1-\rho(k,C)+o(1)$. Then (T) fails.
\end{corollary}

\begin{outline}
The simulators of a machine $M$ with $k$ work tapes have $k$ work tapes (\lean{native_machine_k}), so for every $c>0$ the deficit
$\rho(k,11c+1)$ satisfies the hypothesis of \cref{thm:G} for $M$ and $c$. At $\rho=\frac12$ the hypothesis is the native half bound of the
fourteenth paper~\cite[Corollary~L]{BoundedStack2026}, assumed uniformly over the machines with $k$ work tapes and over all depths, as in
the eighth paper~\cite[\S6.2]{EighthPaper2026}.
\end{outline}

The fifteenth paper reaches the same conclusion from the half bound on chosen simulators~\cite{HighRate2026}; Theorem~G accepts any
fixed deficit, chosen for each machine and time bound separately and needed only on that machine's simulators. In the uniform form of
\cref{cor:G-uniform} the hypothesis is at least as strong as a uniform crux; its value is that the deficit may be tiny, which is the
form that Theorems~B and~J produce, but only for their classes.

For a single machine of low adaptivity the same counting needs no simulator: a deficit above the entropy of the trace suffices.

\begin{corollary}[A native gap above the trace entropy]\label{cor:G-low}
\formal{NativeGapCriterion.lean}{fSecond_native_gap, native_gap_exclusion_of_trace_support}
Let $M$ be a machine with $k$ work tapes, let $c>0$ and $\eta<\rho$, let $\beta_1$ be the block width of $M$ at scale~$1$
(\lean{MergeDepth.blockWidth}), and put $D=\mathrm{genieDepth}\bigl(k,\beta_1,\max(c,1),(\rho-\eta)/2\bigr)$. Suppose that every binary storage code on
the cone relation with arbitrary roots, after a genie of full depth $D$, of every run of $M$ that halts within $c(L+1)$ steps has rate at most
$1-\rho+o(1)$. Then at all large $d$ at which $M$ transposes within time $c$, the skeleton of $M$ at scale~$1$, and hence its trace, has
entropy greater than $\eta L$ on a uniformly random input (\lean{TapeStorageRates.skeletonOfTrace_eq}, \lean{MergeDepth.entropy_comp_le}).
Moreover, if a machine $N$ satisfies the same hypothesis, with its own $k$, $\beta_1$ and $D$, and for all large $d$ every $d$ at which $M$
transposes within time $c$ is one at which $N$ transposes within time $c$ with at most $2^{\eta L}$ distinct traces, then $M$ transposes
within time $c$ at only finitely many~$d$.
\end{corollary}

\begin{outline}
If the trace entropy were at most $\eta L$, some trace would be followed by at least $2^{(1-\eta)L}$ inputs. After the genie of cost
$\frac12(\rho-\eta)L$, its fibre splits into storage codes on the cone relation, one of which has rate at least
$1-\eta-\frac12(\rho-\eta)>1-\rho$, contradicting the hypothesis. This is the argument of \cref{thm:G} with $M$ in place of the
simulator; at $\rho=\frac12$ it is the fourteenth paper's low-adaptivity corollary~\cite[Corollary~L]{BoundedStack2026}.
\end{outline}

\subsection{Few scales and no free head}

For a work head, a \emph{family} is the set of its revisits on one side that share their \emph{turning point}, the first step inside the
revisit at which the head is farthest from the cell; its \emph{height} is the distance from the turning point to the family's outermost
cell, the cell of its longest revisit, as in the eighth paper~\cite{EighthPaper2026}. The head has \emph{owner load} at most $\ell$ if
at every step, on each side, the revisits open at that step have at most $\ell$ distinct turning points. Bounded cell entry $V$ gives owner
load at most $2V-1$.

\begin{theorem}[Theorem~Q, scale-gap pigeonhole]\label{thm:Q}
\formal{ZeroFreeHeadQWindow.lean}{zeroFreeHead_qWindow_excludes, zeroFreeHead_entries_qWindow_excludes, zeroFreeHead_qWindow_oblivious_excludes}
Let $c,\beta$ be reals and $k,\ell,q,h$ natural numbers. There are $\eps>0$ and $d_0$ such that for every $\eta<\eps$ and $d\ge d_0$ no machine with at
most $k$ work tapes and value width at most $\beta$ transposes at size $d$ within time $c$ with an $\eta$-heavy trace on which every work
head has owner load at most $\ell$ and the binary logarithms of the heights of all families lie in at most $q$ windows of width~$h$. The
input head is unrestricted. For oblivious machines the statement holds at every $d\ge d_0$ with $\eta=0$.
\end{theorem}

\begin{outline}
Owner load $\ell$ colours the families with $N=2k\max(\ell,1)$ classes of non-overlapping spans, so their heights sum to less than $NcL$. Choose
thresholds $Q_0<Q_1<\dots<Q_q$, each so large compared with the last that a window can be \emph{bad}, that is contain the height of a
family longer than the current threshold, for at most one index. With $q$ windows and $q+1$ indices some index is bad for no window, so
only boundedly many families are large, and the eleventh paper's low-adaptivity theorem for boundedly many large
families~\cite{MergeDepth2026} gives the contradiction. Its constant comes from a compactness argument and is not
explicit~\cite{SourceCovers2026}.
\end{outline}

Theorem~Q is the case without a free work head of the eighth paper's bounded-load, boundedly-many-scales class, which is
conditional on the crux there~\cite{EighthPaper2026}; here it is unconditional, with any input motion and low adaptivity. With
$\Theta(\log d)$ scales instead of boundedly many, the question is open already for two work tapes.

\section{What a linear-time transposer needs}\label{sec:construction}

This section reads the results backwards, as a specification that any oblivious linear-time transposer must meet, and then
describes how close known schedules come to meeting it.

\subsection{The specification}

Suppose $M$ is an oblivious machine with $k$ work tapes that transposes in time $C(L+1)$ at infinitely many sizes. By the
results of this paper and of the series, at all large such sizes:
\begin{enumerate}[leftmargin=2em,label=(\roman*)]
\item \emph{Shape.} At scale $\theta L/4$ each work tape's walk is one of finitely many words, and some tape turns at least
twice with overlapping legs, so it reuses a stretch of tape like a queue (\cref{thm:S}). Along infinitely many sizes the shape
is one fixed tuple of words.
\item \emph{Crossings.} From every set of at least $\alpha L$ outputs, a linear number of vertex-disjoint dependency paths reach
the earlier ends of crossing revisits of excursion more than $\theta L$ (\cref{thm:X}), and the crossing steps of one coarsely monotone leg
of one tape already carry a linear fraction of them; the tape and the leg may depend on the output set and on the size, and are fixed along
infinitely many sizes for the set of all outputs (\cref{thm:S}).
\item \emph{Two long-range heads.} At least two work heads revisit cells after moving a fixed fraction of $L$ away, and for each work head
some set containing a fixed fraction of the outputs reaches long revisits off that head with linear expansion (\cref{thm:H}).
\item \emph{Rows to columns.} For every two sets of rows $A_0,A_1$, every set of steps separating the reads of rows $A_0$ from
the outputs that rows $A_0$ feed into output rows $A_1$ has at least $\abs{A_0}\abs{A_1}b/(\wval+1)$ elements (\cref{lem:R}).
\item \emph{Codes.} After a genie of cost $\eps L$ and bounded depth, the cone relations of the schedule carry an index code of at most
$\eps L$ bits, one of whose fibres is a storage code of rate at least $1-\eps$~\cite[Corollary~O]{BoundedStack2026}; by (i)--(iii), the
relations involved are those of macroscopic queue-like reuse.
\end{enumerate}
Conditions (i)--(iv) are proved here without any coding hypothesis. Condition (v) is the coding question itself: a linear-time
oblivious transposer is a code of rate close to $1$ on the cone relations of a schedule meeting (i)--(iv). Machine-free relations carry
storage codes above rate one half~\cite{StorageCodes2026}, and binary codes of rate approaching one exist on suitable separated
relations~\cite{HighRate2026}.

\subsection{Mechanisms that fail}

Each of the following design patterns is excluded for oblivious linear-time transposition, whatever the machine computes on
its tapes.
\begin{itemize}[leftmargin=1.5em]
\item \emph{Stack-like reuse at macroscopic scale.} A schedule whose long revisits are nested, so that every tape reuses its
storage like a stack at scale $\theta L$, has no macroscopic crossing steps, and is excluded by \cref{thm:X}.
\item \emph{One long-range head.} One unrestricted work head together with heads whose revisits stay below a small fixed fraction
of $L$ is excluded by \cref{thm:H}, whatever the input head does.
\item \emph{Boundedly many macroscopic phases.} Schedules in which the input head and some bank heads are weakly monotone on boundedly many
phases while the other heads move inside sublinear workspaces are excluded by \cref{thm:BM}; with sublinear jitter on the bank heads, and
for adaptive machines, by \cref{thm:J}.
\item \emph{Cheap row-to-column cuts.} A schedule with a set of fewer than $\abs{A_0}\abs{A_1}b/(\wval+1)$ steps separating the reads
of some rows $A_0$ from their outputs in output rows $A_1$ is excluded by \cref{lem:R}; \cref{thm:colouring} gives a convenient
way to exhibit such cuts.
\item \emph{The explicit families of \cref{app:catalogue}}, each excluded by a class theorem of \cref{sec:classes} or a certificate of
\cref{sec:certificates} (in proof sketches). Several of them make macroscopic crossings with linear expansion (numerically, at $d\le128$), so
\cref{thm:X} alone does not exclude them.
\end{itemize}

\subsection{Schedules that meet every condition we could test}\label{sec:cascades}

Five explicit oblivious schedules, which we call \emph{two-sided cascades}, meet every necessary condition of this paper that we could
test. Each loads the input onto two or three work tapes and then runs a cascade over the scales $M_r=L/2^r$, $1\le r<\log_2d$. A forward
stage folds the data at each scale: one half by a transfer at varying relative speed, the other merged in with a cyclic shift of about
$M_r/3$. A backward stage unfolds it by $1$-to-$2$ doublings with an unshifted and a shifted half. A final sweep emits the $L$ output bits.
Every stage at scale $r$ moves the heads over $O(M_r)$ cells a bounded number of times, so the running time is $O(\sum_rM_r)=O(L)$. The five
differ as follows.
\begin{enumerate}[leftmargin=2em]
\item Two work tapes; before the cascade, two symmetric pairs of two-dimensional cyclic shifts of the entry grid are applied to a copy of
the input.
\item Two work tapes; before the cascade, an anti-diagonal reflection and two butterfly shifts, by $2^{m-1}(d-1)b$ and $(d-1)b$ where
$d=2^m$, are applied to a copy.
\item Three work tapes, every cell entered at most ten times; the third tape holds disjoint zones for the folded data at all scales and
moves through them monotonically in each stage.
\item Three work tapes; the forward and backward stages use disjoint zones for all scales, and no zone is overwritten.
\item Two work tapes; the folds are $2$-to-$1$ compressions, and every scale has a comb pass in which one tape makes $2^r$ teeth while the
other waits.
\end{enumerate}
The machines that realise them in our experiments compute other functions; whether some machine with one of these oblivious traces
transposes is \cref{prob:cascades}.

\begin{remark}[Numerical]\label{rem:cascades}
At $d\le128$ and $b=1$ the measured running times of the five schedules lay between $18L$ and $90L$ and stabilised as $d$ grew. On all five,
the minimum set separating a large output set from the earlier ends of crossing revisits of excursion greater than $E$ had exactly as many
elements as the output set, for every tested $E\le L/2$. Their skeletons at scale about $L/6$ were independent of $d$, with a single leg
carrying linear expansion, and their numbers of disjoint expanding sets of steps stayed bounded, as linear time requires. So, as far as
these measurements go, the conditions (i)--(iv) of the specification do not distinguish these schedules from a transposer. A one-pass shear,
excluded by \cref{thm:J}, and a multi-scale stack schedule, excluded by a certificate (in a proof sketch, \cref{app:catalogue}), also make
macroscopic crossings with linear expansion in these measurements, which suggests that (i)--(iv) are not sufficient for transposition.
\end{remark}


\subsection{Why cuts alone cannot finish}

The certificates of this paper are cuts: they bound from below the number of steps that must separate inputs from outputs. A
single cut can only certify storage, and the next proposition shows that for transposition one set of steps serves as the cut
for every pair of row sets at once.

\begin{proposition}[One cut for every pair of row sets]\label{prop:onecut}
\notformalized
Let a run with horizon $T$ read its last input cell at step $t_{\mathrm{in}}$, and put $t^*=t_{\mathrm{in}}+1$. Let $W^*$ consist of
\begin{itemize}[leftmargin=1.5em]
\item every step $u<t^*$ that is the earlier end of a revisit $(u,t)$ with $t\ge t^*$,
\item every output step before $t^*$, and
\item the step $t^*$, if $t^*<T$.
\end{itemize}
Then for every two sets $A_0,A_1$ of rows, $W^*$ meets every dependency path from the output steps of the entries that rows $A_0$
feed into output rows $A_1$ to the steps that read rows $A_0$. The steps of the first kind number at most the number of work cells that
are visited both before $t^*$ and at or after~$t^*$.
\end{proposition}

\begin{informalproof}
Take a dependency path from an output step $t_0$ to an input-reading step $u$. Every edge of the dependency DAG goes from a step to
a strictly earlier step: the state edge $a\to a-1$, or the revisit edge $a\to a'$ for a revisit $(a',a)$. If $t_0<t^*$, then $t_0$ is
an output step before $t^*$ and lies in $W^*$. Otherwise $t_0\ge t^*>t_{\mathrm{in}}\ge u$, so the path has an edge $a\to a'$ with
$a\ge t^*>a'$. If it is a state edge, then $a'=a-1$, so $a=t^*$, and $t^*\in W^*$. If it is a revisit edge, then $(a',a)$ is a revisit
with $a'<t^*\le a$, so $a'\in W^*$. In each case the path meets $W^*$; this argument does not use $A_0$ or $A_1$. A step of the first kind
is the earlier end of a revisit of some tape that spans $t^*$; on each tape, each cell is the cell of at most one such revisit, namely
the one from its last visit before $t^*$ to its next visit.
\end{informalproof}

For a transposer, \cref{lem:R} with $A_0=A_1$ the set of all rows gives $\abs{W^*}\ge(1-\eta)L/(\wval+1)$. When few outputs precede the last
input read, this says that the storage live across $t^*$ is linear. In any case \cref{prop:onecut} shows that this one bound already
accounts for every pair of row sets: however finely the rows are split, cut certificates for different splits may coincide, and a lower bound
beyond linear time cannot come from counting disjoint cuts unless something other than the cut property forces the cuts apart. The
colouring criterion of \cref{sec:demand} verifies such a separation on specific schedules (\cref{cor:scale-separated}), but it does not supply
one in general.

Outside such restricted classes, the only coding-proof logarithmic factor known on tapes is of a different kind. For an oblivious transposer
whose run has $r$ monotone passes, Theorem~P of the first paper gives $k\,db\le((k+1)^r-1)(2b+\sigma_M)$, where $\sigma_M$ is the number of
bits of a state and the scanned symbols, so a fixed machine needs $r\ge\log_{k+1}d-O(1)$ passes~\cite[Theorem~5.9]{FirstPaper2026}; for
adaptive machines an extra term appears~\cite[Theorem~5.12]{FirstPaper2026}. Its grouped form~\cite[Theorem~5.8]{FirstPaper2026} reads
$d^2b\le H_{\mathrm{addr}}\bigl((\gamma+1)db+\lceil d/\gamma\rceil\sigma_M\bigr)$ for every $\gamma\ge1$, where $H_{\mathrm{addr}}=((k+1)^r-1)/k$
bounds the number of addresses of the forests in its proof.

\begin{remark}[Pass counting at one scale]\label{rem:onescale}
Proof sketch only: suppose a run splits into $r$ phases in each of which every head stays within distance $S$ of a monotone trend. Adapting
the proof of the grouped bound so that the nodes of its forests carry the windows of $O(S)$ cells around the trends, at the cost of
bookkeeping for the offsets within each phase, plays the role of a finite control of $\sigma\approx\sigma_M+O(kS\log(g+1))$ bits, and choosing
$\gamma\approx\sqrt{\sigma/b}$ in the grouped inequality gives $(k+1)^r\gtrsim\tfrac k2\sqrt{L/\sigma}$. So boundedly many phases with $S=o(L)$
cannot transpose. At $S=\eps L$, however, the bound only asks for $r\ge\frac12\log_{k+1}(1/\eps)-O(1)$ phases. A linear-time run may have up to
about $4kC/\eps$ legs at scale $\eps L$ (\cref{lem:finite}), and schedules whose number of legs at scale $\eps L$ grows like $\log(1/\eps)$, such
as the cascades of \cref{sec:cascades}, meet the bound at every scale.
\end{remark}

So pass counting, like cut counting, cannot exclude linear time one scale at a time. \Cref{prob:multiscale} asks for a pass bound
that counts phases across all scales at once.

\section{Open problems}\label{sec:open}

The results of this paper confine an oblivious linear-time transposer to macroscopic queue-like reuse with linear expansion, on one
leg of a bounded skeleton, with two long-range heads. The following problems would take the next step; the first three are on the
side of impossibility, the fourth on both sides.

\begin{problem}[Pass bounds across scales]\label{prob:multiscale}
For a run and $1\le j\le\log_2L$ let $r_j$ be the total number of legs of its skeletons at scale $\lfloor2^{-j}L\rfloor$. A linear-time run has
$r_j\le k(4C\,2^j+1)$ (\cref{lem:finite}), and pass counting at a single scale requires only $r_j\gtrsim\frac j2\log_{k+1}2$
(\cref{rem:onescale}). Is there, for every $k$ and $C$, a function $f$ with $f(d)\to\infty$ such that no oblivious machine with $k$ work tapes
whose runs satisfy $r_j\le Cj$ for all $j$ transposes at size $d$ in time less than $f(d)L$? More generally, find a coding-proof inequality
that involves all the $r_j$ at once and fails for every linear-time run.
\end{problem}

The schedules with $r_j=O(j)$ include, as far as we can tell, the cascades of \cref{sec:cascades}. A positive answer would exclude every
transposer whose motion looks like a cascade at all scales; a negative answer would be a transposer of that shape, faster than
$O(L\log d)$.

\begin{problem}[From the simulator back to the machine]\label{prob:pullback}
Let $M$ transpose in linear time, and let $\Ns$ be its bit simulator. Does the conclusion of \cref{thm:X} hold in $M$'s own dependency
DAG, on a set of traces followed by all but an exponentially small fraction of the inputs? A positive answer would remove the word
``oblivious'' from Theorem~X. The question is whether a macroscopic crossing in the dependency DAG of $\Ns$, which rewrites windows of $3s$
cells in each round, forces a macroscopic crossing of the heads of~$M$.
\end{problem}

\begin{problem}[Constants]\label{prob:constants}
Does \cref{thm:X} hold with a threshold $\theta\ge c_1/C$ and an expansion constant $c_2$, where $c_1,c_2>0$ depend polynomially on $1/k$,
$1/\wval$ and $\alpha$? Numerically, the threshold of Theorem~D appears not to be a universal constant (\cref{rem:numerics-crossings}), but
nothing we know rules out thresholds of order $1/C$.
\end{problem}

\begin{problem}[The cascades]\label{prob:cascades}
For each of the five cascade schedules of \cref{sec:cascades}, decide whether some machine with that oblivious trace transposes. Each is a
linear-time schedule that met every necessary condition of this paper in our tests (\cref{rem:cascades}). A transposer with one of these
traces would prove~(T); an exclusion of one of them would, as far as our tests indicate, need an argument that uses more than
Theorems~D, X, S and Lemma~R.
\end{problem}

A linear-time transposer, if one exists, must solve a coding problem on its own schedule: codes of rate approaching one on the cone
relations of macroscopic queue-like reuse~\cite[Corollary~O]{BoundedStack2026}. The fifteenth paper shows that binary storage codes of
rate approaching one exist on suitable separated relations~\cite{HighRate2026}; realizing such codes on the relations of a schedule that
meets conditions (i)--(iv) of \cref{sec:construction} would be a construction. In the other direction, showing that the cone relations of
every such schedule, after a cheap genie, have a fixed rate deficit would exclude oblivious linear-time transposers
(\cref{thm:crux-excludes};~\cite[Corollary~O]{BoundedStack2026}); deciding (T) itself needs the uniform native form of \cref{thm:G}.


\appendix
\section{Catalogue of excluded schedule families}\label{app:catalogue}

This appendix lists the explicit oblivious linear-time schedule families excluded by the results of this paper. A schedule family
prescribes, for each size $d$ in its domain, the motion of the input head and of the work heads and the append steps; a family is
\emph{excluded} if for every $C\ge1$ and every value width $w$ there is $d_0$ such that, at every size $d\ge d_0$ in its domain, no machine of
value width at most $w$ whose trace is the prescribed one transposes within $C(L+1)$ steps.

\paragraph{Status.} For rows 1--20 the excluding theorem is formal, and that the schedule belongs to the class is checked by inspection
of the schedule. \Cref{cor:H} and \cref{thm:Q} exclude size by size, with thresholds that depend only on the bounds for the number of tapes and the
value width (uniform over all machines within those bounds),
so the rows they cover are excluded in the sense above. \Cref{thm:BM,thm:J} speak of one machine at all large sizes: for the rows that
only they cover (rows 3--7 and 12) they exclude the machines that follow the schedule at all large sizes of its domain, not, size by size,
every machine that follows it at one size. For the schedules defined only for even $d$ or for powers of two (rows 3, 5 and 6) they are
applied along those sizes, which their formal class predicates, quantified over all large $d$, do not literally cover. For rows 21--38 the certificate is
built from formal lemmas of \cref{sec:certificates}, but its verification on the schedule is a written argument; each is stated below as a
proof sketch only. \emph{None of these exclusions is a numbered result of this paper.}

\paragraph{Crossings (numerical).} The last column records whether \cref{thm:X} alone could exclude the family. ``Yes'' means that the
schedule's work heads stay within $O(b)$ cells, so it has no crossing revisits of linear excursion at all. The other entries are
numerical, from exact computations at $d\le128$: ``no'' when the minimum set separating a large output set from the crossing steps of
revisits of excursion more than $2W$ had at least $0.25$ times as many elements as the output set; ``small'' when that ratio was at most
$0.07$, which suggests, as a heuristic, that \cref{thm:X} applies; and ``---'' when nothing was measured.

{\small
\begin{longtable}{@{}r p{0.42\textwidth} c p{0.2\textwidth} c@{}}
\toprule
& Schedule & $k$ & Excluded by & Crossings \\
\midrule
\endhead
1 & Direct copy: output bit $\ell$ emitted while reading input bit $\ell$; work tape idle (computes the identity) & 1 & \cref{cor:H}, \cref{thm:BM} & yes \\
2 & Persistent triangles: output $(j,i,s)$ emitted while reading input $(j,-j-i\bmod d,s)$; work head within $b$ cells & 1 & \cref{cor:H}, \cref{thm:J} & yes \\
3 & One-pass shear with row rotations, two banks and a row buffer; seven monotone phases ($d$ even) & 3 & \cref{thm:BM,thm:J} & no \\
4 & Ruler schedule: two banks moving one row per episode, a third head making excursions of $2$-adic lengths inside one row & 3 & \cref{thm:J} & --- \\
5 & Row expander: rows mixed in a two-head deep-nesting workspace of width $W$, two banks, four monotone phases ($d$ even) & 4 & \cref{thm:BM} & --- \\
6 & Fan-out cascade fed by fresh input slices at each scale ($d=2^m$) & 2 & \cref{thm:J} (suffix form) & --- \\
7 & Unit-slope passes: load, reflection, shifted sweep, emit; at most three constant-velocity segments per pass & 2 & \cref{thm:J} & small \\
8 & Bounce: zig-zag input, one tape sweeping $[0,L)$, nested excursions of height at most $L/(4d)$ on the other & 2 & \cref{thm:J}, \cref{cor:H} & small \\
9 & Sublinear stack: blocks of $L/(4d)$ outputs pushed and popped on fresh cells while the other tape sweeps & 2 & \cref{thm:J}, \cref{cor:H} & --- \\
10--11 & Free pseudo-random Dyck walk of length $L$ on one tape while the other sweeps the tape six times (two variants) & 2 & \cref{thm:J}, \cref{cor:H} & --- \\
12 & One push and one pop of height $2L$ on a stack tape while the other tape bounces & 2 & \cref{thm:J} & --- \\
13--16 & One free head with all other excursions $O(L/\log d)$: slot cascades and skewed wavefronts (four variants) & 2 & \cref{cor:H} & small (16) \\
17--20 & No free head, owner load at most $4$, family heights in at most $4$ scale windows: multi-scale push--pop, two-excursion hub, LIFO Dyck stack, torus shifts & 2 & \cref{thm:Q} & --- \\
21--23 & One push--pop excursion on a stack tape, with a free tape running a multi-scale cascade or bouncing in a hub (three variants) & 2 & slab-chunk genie on a two-scale grid & --- \\
24 & Multi-scale Dyck stack: $d$ sawtooth waves on a stack tape while the free tape runs a cascade and bounces in a hub ($d=2^\kappa$) & 2 & block cut on the hub & no \\
25--26 & Two push--pop excursions on a stack tape, with a mandatory cut at the second excursion (two variants) & 2 & block cut & --- \\
27 & Two-exception inverse-square geometry, swept variant~\cite{GapAreaMemo2026} ($d\ge64$) & 10 & closed post-region, $O(d^{1/3})$ boundary steps & no \\
28--29 & Two-exception inverse-square geometries~\cite{GapAreaMemo2026} ($d\ge64$) & 9, 11 & window lemma after two cut steps & small \\
30 & Mirrored dyadic fold cascade & 2 & window lemma (\cref{lem:stepinv}), uncut & small \\
31 & Unit-slope dyadic cascade ($d=2^m$) & 2 & sub-row window, uncut & --- \\
32 & Dyadic fold cascade with row blocks of doubling width & 2 & window lemma, uncut & no \\
33--34 & Two stack excursions with the free tape in a hub of width $L/2^s$ ($s$ fixed) or $d$ & 2 & block cut; space--time potential & --- \\
35 & Daisy-chained stack copies with golden-ratio shifts & 2 & piecewise-isometric windows & --- \\
36 & LIFO fan-out cascades of $\log d$ levels (two variants; $d=2^\kappa$) & 2 & laminar cut of $2^{1-p}L+3$ steps, $p=O(\log\wval)$ & no \\
37 & Equal-volume tapped cascade with $\log d$ meso-scale stages & 2 & meso-hub cut of $O(L/\log d)$ steps & no \\
38 & Two-band meso cascade & 2 & closed post-region, $O(L/\log d)$ boundary steps & --- \\
\bottomrule
\end{longtable}
}

\paragraph{Rows 1--20.} For bounded macro-phases (\cref{thm:BM,thm:J}) the crux rate is $1/(4(\abs B+1)^{M_{\mathrm{ph}}-1})$ with the bank $B$ and
the number of phases read off the schedule: for example, row~3 has banks $\{0,1\}$ and seven phases, row~5 banks $\{0,1\}$ and four phases.
Rows 13--16 satisfy the hypotheses of \cref{cor:H} with the free head named and all other excursions $o(L)$; rows 17--20 those of
\cref{thm:Q} with owner load and window counts read off the schedule. For row~1, \cref{thm:BM} gives crux rate $\frac14$; the true rate is at
most $\frac12$, because pairing $q$ with $\tpos(q)$ gives an exclusive-or code of $L/2$ bits.

\paragraph{Rows 21--38.} Proof sketch only. Rows 21--23 are excluded by the frozen-trace count used in the proof of \cref{thm:H}: on the pop
leg of the stack, revisits that span the peak are order-reversing, and a two-scale slab grid with a greedy slicing gives a trapped chunk whose
genie costs less than its output surplus. Rows 24--38 are excluded through \cref{thm:crux-excludes} by an explicit certificate. Rows 24--26 and 33:
the cut at a mandatory part of the stack motion makes every relevant output carry at most $a$ blocks of potential sources and every input
at most $e$ blocks, and \cref{lem:blockcut} gives an acyclic set of density $1/\binom{a+e}{a}$, uniformly in the depth. Rows 27--30 and 32: the
window lemma (\cref{lem:window,lem:stepinv}) with potentials $\varphi(q)=\tpos(q)$ and $\psi(p)=L-1-p$ or a row-local variant, possibly after
cutting a sublinear boundary. The geometries of rows 27--29 have hop counts of order $L/i^2$ in their $i$-th family; the gap-area memo of the
ninth paper shows that its frozen-gap-area certificate cannot exclude this case~\cite{GapAreaMemo2026}, whereas the certificates here use cut
steps and a window potential instead. Row 34 uses a space--time potential, row 35 a bound for piecewise-isometric motion with boundedly many
segments, row 36 a bounded-level laminar cut, row 37 a meso-scale hub cut, and row 38 the formal closed-post-region criterion
(\lean{isAcyclicIn_ballAvoidSystem_of_closedPostRegion}).

As a numerical check on the method, the same classification was run on explicit transposers that run in time $\omega(L)$ (shuffle and
butterfly transposers, a radix transposer on a stack, and one-tape transposers); none of the exclusions was triggered by them.

\section{Index of Lean declarations}\label{app:leanindex}

% Draft placeholder. This appendix is generated from the LaTeX sources, together with lean/PaperAxioms.lean, when the Lean
% library is added (as in the fourth paper's tools/lean_index.py). Until then it only describes the library.

\begin{sloppypar}
The formal proofs are in the Lean~4 library \leanfile{NativeRelations} in \leanfile{lean/} (Lean and Mathlib v4.34.1). It requires the
libraries of the first, eighth, eleventh, fourteenth and fifteenth papers of the series by relative path. Its modules are, in the order of
this paper: \leanfile{DichromaticCrux.lean}, \leanfile{MachineCrux.lean} and \leanfile{EnvelopeCrux.lean} (\cref{sec:certificates});
\leanfile{OneFreeHeadSmallExcursion.lean} (\cref{sec:heads}); \leanfile{ConjectureEBridge.lean}, \leanfile{MacroscopicThreshold.lean}
and \leanfile{TightConstants.lean} (\cref{sec:crossings}; the last holds the forms with the constants of \cref{sec:constants}, which
also serve \cref{sec:skeletons,sec:adaptive}, and \cref{lem:acyclic-cost}(b), and the other two the first forms); \leanfile{MacroscopicSkeleton.lean} (\cref{sec:skeletons}); \leanfile{TranspositionCut.lean} (\cref{sec:demand});
\leanfile{SimulatorBridge.lean} (\cref{sec:adaptive}); \leanfile{MacroBankCrux.lean}, \leanfile{SuffixMassCrux.lean},
\leanfile{JitterMacroBank.lean}, \leanfile{AdaptiveSuffixJitterBM.lean}, \leanfile{NativeGapCriterion.lean} and
\leanfile{ZeroFreeHeadQWindow.lean} (\cref{sec:classes}); \leanfile{CoarseGroupedPass.lean} (\cref{sec:construction}); and \leanfile{ScaleProfile.lean} (\cref{rem:waves}). The complete index of the
declarations cited in this paper, with their files, will be generated from the sources.
\end{sloppypar}

\section{Building and verifying the formalization}\label{app:verification}

% Draft placeholder, to be completed when the Lean library is added.

The library builds with the Lean toolchain and the Mathlib revision pinned in \leanfile{lean/lean-toolchain} and
\leanfile{lean/lake-manifest.json}, the same as the earlier papers of the series. Because it requires sibling libraries by relative path,
the whole repository must be checked out. From \leanfile{lean/}:
\begin{verbatim}
lake exe cache get
lake build
lake env lean PaperAxioms.lean
\end{verbatim}
The last command prints, for every declaration cited in this paper, the axioms it depends on; each depends at most on
\lean{propext}, \lean{Classical.choice} and \lean{Quot.sound}. The library contains no \texttt{sorry}, no \texttt{native\_decide}, no new
axioms and no code that bypasses the kernel.

\section{Proofs for Cascades 3 and 4}\label{app:cascade-cuts}

This appendix proves \cref{thm:cascade-cuts} and \cref{rem:cascade-widths}. Throughout, $d=2^k\ge32$, $b\ge1$, $L=d^2b$,
$m_r=L/2^r$, $\mathrm{sh}_r=\lfloor m_r/3\rfloor+b$ and $r_{\max}=k-1$, so that $m_r\ge m_{r_{\max}}=2db\ge64$ for $r\le r_{\max}$. Let
$A=[db,3db)$ be the set of input positions of rows~$1$ and~$2$. As $A_1$ is the set of all rows, $\Qstar(\{1,2\},A_1)=A$, and
$\Lin(\{1,2\})$ is the set of steps at which the input head reads a position of $A$. The traces are those of the constructors
\texttt{build\_c3\_ten\_entry\_cascade} and \texttt{build\_c4\_no\_overwrite\_cascade} in \texttt{experiments/src/schedules.rs}; the head
motions below are read off their loops, and we also checked them on the exported traces for $d\le128$. A \emph{region} is an interval
of cells of one tape, and the \emph{offset} of a cell in a region is its position minus the first position of the region.

\subsection{Taint closures}

For a work tape $j$ and a step $t$, let $\prev_j(t)$ be the last step before $t$ at which head $j$ scans the cell it scans at step $t$,
if there is one; the cell edge into $t$ of tape $j$ is then $\prev_j(t)\to t$. The \emph{wait age} $\omega_j(t)$ is $0$ if $t=0$ or
head $j$ scans different cells at steps $t-1$ and $t$, and $\omega_j(t-1)+1$ otherwise.

\begin{definition}[Taint closure]\label{def:taint}
Fix a set $C_0$ of cell edges and an order $h\in\{1,2\}$. A step $t$ is \emph{tainted} if it reads a position of $A$, or if for some
tape $j$ the cell edge $p\to t$ with $p=\prev_j(t)$ is not in $C_0$, the step $p$ is tainted, and either $p\le t-2$ or $\omega_j(t)\le h$.
A \emph{run} is a maximal interval of consecutive tainted steps. The \emph{cut} $C$ consists of $C_0$ and of the state edge and the
waits that leave the last step of each run. A cell of tape $j$ is \emph{live tainted} at step $t$ if its last visit $p<t$ is tainted and
the edge $p\to\mathrm{next}_j(p)$ is not in $C_0$.
\end{definition}

A step $t$ with $p=\prev_j(t)=t-1$ is a \emph{short wait} if $\omega_j(t)\le h$. For $h=1$ this means that head $j$ arrived on the cell at
step $t-1$.

\begin{lemma}[The closure is cut by $C$]\label{lem:closure}
\notformalized
Every edge of the cell network from a tainted step to an untainted step lies in $C$. Let $N_{\mathrm{run}}$ be the number of runs and
$N_{\mathrm{up}}$ the number of tainted steps of $\Out(A)$. Then:
\begin{enumerate}[label=(\alph*)]
\item every path in the cell network from a step that reads a position of $A$ to an untainted step uses an edge of $C$;
\item the set $\Gamma$ of the tails of the edges of $C$ and the tainted steps of $\Out(A)$ separates $\Out(A)$ from $\Lin(\{1,2\})$
in the dependency DAG;
\item $C$ has at most $\abs{C_0}+3N_{\mathrm{run}}$ cell edges and at most $N_{\mathrm{run}}$ state edges, and
$\abs\Gamma\le\abs{C_0}+N_{\mathrm{run}}+N_{\mathrm{up}}$.
\end{enumerate}
\end{lemma}

\begin{informalproof}
Let $p\to v$ be an edge with $p$ tainted and $v$ untainted. A state edge has $v=p+1$, so $p$ ends a run and the edge is in $C$. A cell
edge with $v\ge p+2$ is in $C_0$, since otherwise $v$ would be tainted. A wait of tape $j$ that is not in $C_0$ has $\omega_j(v)>h$,
since otherwise $v$ would be tainted; so $p$ ends a run, and the wait is in $C$. This proves the first claim, and (a): a path from a
tainted step to an untainted one has a first untainted vertex, and the edge into it leaves a tainted vertex.

For (b), take a dependency path from $\Out(v)$, $v\in A$, to a step that reads a position of $A$. If $\Out(v)$ is tainted it lies in
$\Gamma$. Otherwise the reversed path is a path in the cell network (\cref{def:cellnet}) from a tainted step to an untainted one, and
the edge into its first untainted vertex lies in $C$ and leaves a vertex of the path, which is a tail of $C$. For (c), $C_0$ has
$\abs{C_0}$ edges and tails, and every other edge of $C$ leaves the last step of a run: one state edge, none after the last step of
the trace, and at most one wait per work tape; both schedules have three.
\end{informalproof}

So it suffices to choose $C_0$ so that $\abs{C_0}$, $N_{\mathrm{run}}$ and $N_{\mathrm{up}}$ are small.

\subsection{Cascade 3}\label{app:cascade3}

Cascade~3 uses the closure of order $h=1$. Its regions are as follows. On tape~$0$, $R_r$ starts at $a_r$, with $a_1=L$ and
$a_{r+1}=a_r+2m_r$; for the backward stage, $\mathrm{IN}_r$ has size $m_r$ and $\mathrm{OUT}_r$ has size $2m_r$, with
$\mathrm{IN}_{r-1}=\mathrm{OUT}_r$. On tape~$1$, the level-$r$ scratch starts at some $s_r$, and the backward stage uses fresh regions
$\mathrm{CB}_r$ of size $m_r$. On tape~$2$, $Z_r=[z_r,z_r+m_r)$ with $z_r=L-2m_r$. Put $\Delta=\lfloor\lfloor L/3\rfloor/2\rfloor-1$. The
phases that matter are the following.

\begin{center}\small
\begin{longtable}{>{\raggedright\arraybackslash}p{0.2\textwidth}>{\raggedright\arraybackslash}p{0.74\textwidth}}
\toprule
phase & head motion \\
\midrule
\endhead
load, step $u<L$ & input head on $u$; tapes $0$ and $1$ on cell $u$; tape~$2$ on cell $L+\lceil u/2\rceil$ while $u\le\lfloor L/3\rfloor$,
then at speed~$1$ \\
rev & tape~$0$ down $[L,2L)$ and tape~$2$ down to cell $L$, tape-$2$ cell $x$ aligned with tape-$0$ cell $x+\Delta$; then tape~$2$
parks on cell~$L$; tape~$1$ idle \\
fwd $r$.transfer, step $s$ & tape~$0$ on $R_r$ offset $\min(s,m_r-1)$; tape~$1$ on scratch offset $\lceil s/2\rceil$ for
$s<2\lfloor m_r/3\rfloor$, then at speed~$1$; tape~$2$ idle \\
fwd $r$.move2 & tape~$0$ up to $R_r$ offset $2m_r-\mathrm{sh}_r$, then idle; tape~$1$ down the scratch from offset $m_r-1$ to $0$ \\
fwd $r$.copyS & tapes $0,1,2$ on $R_r$ offset $2m_r-\mathrm{sh}_r+u$, scratch offset $u$, $Z_r$ offset $u$, for $0\le u<\mathrm{sh}_r$ \\
fwd $r$.move3, copyR, move4 & tape~$0$ down to $R_r$ offset $m_r$; then $R_r$ offset $m_r+u$, scratch offset $\mathrm{sh}_r+u$, $Z_r$
offset $\mathrm{sh}_r+u$; then tape~$0$ up to $R_{r+1}$ offset $m_r-1$ \\
fwd $r$.back, move5 & tape~$0$ down $R_{r+1}$ and tape~$2$ down $Z_r$, offset aligned with offset; then tape~$2$ up through $Z_r$ to
$z_{r+1}$ \\
comb $r$.move & tape~$0$ up $\mathrm{IN}_r$ from offset~$0$; tape~$1$ up from $\mathrm{CB}_{r+1}$ offset $\mathrm{sh}_{r+1}$ (from the
end of the scratch if $r=r_{\max}$) to $\mathrm{CB}_r$ offset $m_r-1$; $\mathrm{IN}_r$ offset $s$ is aligned with $\mathrm{CB}_{r+1}$
offset $\mathrm{sh}_{r+1}+s$ while $s<m_{r+1}-\mathrm{sh}_{r+1}$, and with $\mathrm{CB}_r$ offset $s-c_r$, where
$c_r=m_{r+1}+2-\mathrm{sh}_{r+1}$, once $s\ge c_r$; then tape~$2$ moves to $z_r+m_r-1$: one step from $z_{r+1}$ if $r<r_{\max}$, and up
through the offsets $0,\dots,m_r-2$ of $Z_{r_{\max}}$ if $r=r_{\max}$ \\
comb $r$.blocks & $2^r$ blocks of $m_r/2^r$ offsets, from the top: all three heads down together over $\mathrm{IN}_r$, $\mathrm{CB}_r$
and $Z_r$ at equal offsets; then tape~$1$ up and down across the block while tapes $0$ and $2$ wait \\
comb $r$.move2 & tape~$0$ up through $\mathrm{IN}_r$ to $\mathrm{OUT}_r$; tapes $1$ and $2$ idle \\
comb $r$.dbl & $\mathrm{CB}_r$ offset $z$ aligned with $\mathrm{OUT}_r$ offsets $2z$ and $2z+1$ \\
comb $r$.move3, shift, move4, rem & tape~$1$ alone down to offset $\mathrm{sh}_r-1$; then $\mathrm{OUT}_r$ offset $y$ aligned with
$\mathrm{CB}_r$ offset $(\lfloor y/2\rfloor+\mathrm{sh}_r)\bmod m_r$, both moving down, with tape~$1$ alone moving up between the two parts \\
rewind & tape~$1$ down from $\mathrm{CB}_1$ offset $\mathrm{sh}_1$ to cell~$0$; tapes $0$ and $2$ idle \\
emit, step $u$ & tape~$0$ on $\mathrm{OUT}_1$ offset $u$; tape~$1$ on cell $u$; tape~$2$ on cell $\min(u,\sum_rm_r-1)$; output bit $u$
appended \\
\bottomrule
\end{longtable}
\end{center}

The positioning phases omitted from the table (the initial positioning, the moves into and out of rev, and the move that starts each
forward level) read no live tainted cell. We use three properties of the trace.
\begin{enumerate}[label=(\roman*)]
\item Every phase is a \emph{segment}, or a few of them: an interval of steps in which every head is either idle throughout, or
monotone and waiting at most one step at a time. To get segments, cut rev, transfer, move2, comb move and emit where a head stops early,
cut each block into its sweeps, and make the last step of copyS, at which all three heads wait, a segment of its own.
\item After the current level, the regions of tape~$0$ below the current one, the scratch regions, and $\mathrm{CB}_{r'}$ for
$r'>r+1$ are not visited again, except by tape~$1$ in the rewind. $Z_r$ is visited after move5 only in comb $r$.blocks and in the emit.
The exceptions are single idle cells: offset~$0$ of $Z_r$ while comb $r$.move2 to rem run and again in comb $(r-1)$.move, and offset~$0$
of $Z_1$ and of $\mathrm{OUT}_1$ during the rewind. Each adds at most one run of at most two steps.
\item In every segment at most one head moves at half speed (in load, transfer, dbl, shift and rem).
\end{enumerate}

\begin{lemma}[Segments]\label{lem:c3-segments}
\notformalized
In a segment of Cascade~3, every tainted step is (1) a step at which a head that is not idle scans a live tainted cell, (2) one of the
first two steps of an idle head on a live tainted cell, (3) a step that reads a position of $A$, or (4) the step right after one of
these, or the step after that when the first is a short wait of an idle head. Hence the number of runs in a segment is at most the
number of maximal intervals of live tainted cells that it reads, over all tapes, plus the number of runs of reads of $A$; and each run
writes one interval of cells on each tape whose head moves.
\end{lemma}

\begin{informalproof}
A monotone head stays on an interval of cells during an interval of steps. Its first step on a cell is a revisit, or at the start of
the segment a wait of either kind, and its second step there, if any, is a short wait. A short wait at the first step $t_0$ of the
segment from a tainted step $t_0-1$ reads a cell that is live tainted at $t_0$, so $t_0$ is of kind~(1). An idle head's first two steps
are a revisit and a short wait, and the rest are waits that are not short. Cells written inside the segment are revisited inside it only
by waits. So a tainted step of none of the kinds (1)--(3) is tainted through a short wait from the tainted step before it. By~(iii), all
short waits of moving heads in the segment belong to one head $j$. After its wait at step $t$, head $j$ moves, so step $t+1$ is no wait
of $j$, and no wait of any other moving head either; so the extension stops after one step. An idle head's short wait at $t_0+1$ can hand
over to a short wait of $j$ at $t_0+2$, so the extension stops after at most two steps. It is adjacent to its run, so it adds no run and
extends each written interval by at most two cells.
\end{informalproof}

The cut is placed at a level $\rho$ with $1\le\rho\le r_{\max}$: $C_0$ is the set of tape-$1$ edges that leave tainted steps of
fwd $\rho$.transfer and whose next visit lies after the transfer. The waits of tape~$1$ inside the transfer are not cut. This is well
defined: no edge of $C_0$ ends at or before the end of the transfer, so the taint up to that point does not depend on $C_0$.

\begin{lemma}[Forward windows]\label{lem:c3-forward}
\notformalized
For $1\le r\le\rho$ let $W_r$ be the set of offsets of $R_r$ that are live tainted when fwd $r$.transfer starts. Then
$W_1=\{\lceil u/2\rceil+\Delta:u\in A\}$ and $W_{r+1}=\{\lceil s/2\rceil:s\in W_r\}$ for $r<\rho$. Each $W_r$ is an interval with
\[
  \abs{W_r}\le\frac{db}{2^{r-1}}+2,\qquad \frac{m_r}3-2<\min W_r\le\max W_r\le\frac{m_r}3+\frac{3db}{2^r}+1 .
\]
Moreover $\abs{C_0}\le db/2^\rho+2$, and the forward stage up to the end of fwd $\rho$.transfer has $O(\rho)$ runs.
\end{lemma}

\begin{informalproof}
The bounds hold for $W_1$: the interval $\{\lceil u/2\rceil:u\in A\}$ has at most $db+1$ elements, from $\lceil db/2\rceil$ to
$\lceil(3db-1)/2\rceil$, and $m_1/3-2\le\Delta\le m_1/3-1$. They pass from $W_r$ to $W_{r+1}$, because $\lceil s/2\rceil\in[s/2,(s+1)/2]$
and the image of an interval of $n$ integers has at most $n/2+1$ elements. Since $3db/2^r=3m_r/d\le3m_r/32$ and $m_r\ge64$,
\begin{equation}\label{eq:c3-window}
  \max W_r\le\tfrac{41}{96}m_r+1<2\lfloor m_r/3\rfloor,\qquad \max\{\lceil s/2\rceil:s\in W_r\}\le\tfrac{41}{192}m_r+1<\mathrm{sh}_r-2 .
\end{equation}

Load: the tainted steps are the reads of $A$ and, at most, the step after the last one, a short wait of tape~$2$. They write tape-$0$
cells, which are not visited again; the loaded copy of rows~$1$ and~$2$ on tape~$1$, cells $[db,3db]$, which is live until the rewind;
and the tape-$2$ cells $L+\lceil u/2\rceil$, $u\in A$, as $3db<\lfloor L/3\rfloor$. Rev reads these tape-$2$ cells and writes the tape-$0$
cells $L+\lceil u/2\rceil+\Delta$, that is, $W_1$; tape~$1$ is idle, so its wait is cut, and the tape-$2$ cells above $L$ are not
visited again.

In fwd $r$.transfer the steps $s\in W_r$ are tainted, and so is $s+1$ for odd $s$, by a short wait of tape~$1$. By \eqref{eq:c3-window}
tape~$1$ is in its half-speed part, and the scratch cells it writes are the offsets $S_r=\{\lceil s/2\rceil:s\in W_r\}$. For $r=\rho$ the
edges leaving them are $C_0$, so $\abs{C_0}\le\abs{S_\rho}\le\abs{W_\rho}/2+1$. For $r<\rho$, the scratch window $S_r$ lies in
$[0,\mathrm{sh}_r-2)$. It is read in move2, where tape~$0$ is already idle because $m_r-1-\max S_r>m_r-\mathrm{sh}_r+1$; in copyS, aligned
with $Z_r$ offset $u$ and with $R_r$ offset $2m_r-\mathrm{sh}_r+u$; and next in the rewind. The copy on $R_r$ is read in move3 and move4
and then never again. The copy on $Z_r$ is read in back, which writes the offsets $S_r$ of $R_{r+1}$, so $W_{r+1}=S_r$, and in move5.
CopyR touches only scratch and $Z_r$ offsets at least $\mathrm{sh}_r$ and $R_r$ offsets below $2m_r-\mathrm{sh}_r$, none of them tainted.
Each of transfer, move2, copyS, move3, move4, back and move5 reads at most two intervals of live tainted cells, so it has at most two
runs (\cref{lem:c3-segments}).
\end{informalproof}

\begin{lemma}[Clean middle]\label{lem:c3-middle}
\notformalized
No step of fwd $r$ for $r>\rho$, of fwd $\rho$ after its transfer, or of comb $r$ for $r\ge\rho$ is tainted.
\end{lemma}

\begin{informalproof}
The cut $C_0$ blocks the scratch window of level $\rho$, so the steps of move2 that read it are untainted and rewrite it with
untainted contents. $Z_\rho$ and $R_{\rho+1}$ are therefore written by untainted steps, and so is everything built from them. These
phases visit only regions of levels at least $\rho$, by~(ii), and the tainted windows of lower levels ($Z_{r'}$, the scratch regions,
$\mathrm{CB}_{r'}$ and the loaded copy) lie in regions they do not visit. The one tainted window of level $\rho$ is
$W_\rho\subseteq R_\rho$; after the transfer, tape~$0$ visits only offsets of $R_\rho$ at least $m_\rho-1$, and $W_\rho\subseteq[0,m_\rho-1)$.
\end{informalproof}

For $r<\rho$ let $U_r$ be the set of offsets $y\in[0,m_r)$ at which $\mathrm{IN}_r$, $Z_r$ or $\mathrm{CB}_r$ is live tainted when
comb $r$.blocks starts, $u_r$ its number of maximal intervals and $\mu_r=\abs{U_r}$; let $h_r=m_r/2^r$.

\begin{lemma}[Backward stage]\label{lem:c3-backward}
\notformalized
We have $u_{\rho-1}\le1$ and $\mu_{\rho-1}\le db/2^{\rho-1}+2$, and for $r<\rho-1$
\[
  u_r\le6u_{r+1}+8,\qquad \mu_r\le10\mu_{r+1}+\frac{db}{2^r}+O(u_{r+1}).
\]
Comb $r$ has at most $c\,(u_r+u_{r+1}+\mu_r/h_r+1)$ runs, for an absolute constant $c$.
\end{lemma}

\begin{informalproof}
Comb $\rho-1$ reads $\mathrm{IN}_{\rho-1}$ and the comb region of level $\rho$, which come from comb $\rho$ and are clean
(\cref{lem:c3-middle}); so $U_{\rho-1}$ is the window $S_{\rho-1}$ of $Z_{\rho-1}$, one interval with at most $db/2^{\rho-1}+2$ cells.

In comb $r$, blocks moves the three heads in alignment, so afterwards all three regions are tainted exactly on $U_r$, together with
one cell at the bottom of a block, where the waiting heads pass taint into the first two steps of the excursion; the excursions only
reread tape-$1$ cells of the same block. Move2 rereads $\mathrm{IN}_r$, which is not visited afterwards. Dbl taints $\mathrm{OUT}_r$ on
$2U_r$; shift and rem, through $y\mapsto(\lfloor y/2\rfloor+\mathrm{sh}_r)\bmod m_r$, add $2(U_r-\mathrm{sh}_r)$ to $\mathrm{OUT}_r$
and $U_r+\mathrm{sh}_r$ to $\mathrm{CB}_r$. So $\mathrm{IN}_{r-1}=\mathrm{OUT}_r$ and $\mathrm{CB}_r$ each have at most $2u_r+2$
intervals, of total size at most $4\mu_r$ and $2\mu_r$.

In comb $(r-1)$.move the alignments of the table add to $\mathrm{IN}_{r-1}$ the image of the comb region of level $r$ (its offsets at
least $\mathrm{sh}_r$, shifted by $\mathrm{sh}_r$), and to $\mathrm{CB}_{r-1}$ the image of the offsets of $\mathrm{IN}_{r-1}$ at least
$c_{r-1}$. The image of $\mathrm{CB}_r$ lands at offsets of $\mathrm{IN}_{r-1}$ below $m_r-\mathrm{sh}_r<c_{r-1}$, where tape~$1$ is still
in the old comb region, so the steps that read it write no cell of $\mathrm{CB}_{r-1}$, and $\mathrm{CB}_{r-1}$ receives at most the
$2u_r+2$ intervals of $\mathrm{OUT}_r$. Hence $u_{r-1}\le(2u_r+2)+(2u_r+2)+2+(2u_r+2)$ and
$\mu_{r-1}\le4\mu_r+2\mu_r+\abs{S_{r-1}}+4\mu_r$, up to $O(1)$ per interval.

Runs: each segment of comb $r$ has at most as many runs as intervals it reads (\cref{lem:c3-segments}). In blocks an interval is read
in every block it meets, and the $u_r$ intervals meet at most $\mu_r/h_r+2u_r$ blocks, three sweeps each. Comb $r$.move reads
$\mathrm{OUT}_{r+1}$ and $\mathrm{CB}_{r+1}$, which gives the term $u_{r+1}$.
\end{informalproof}

Solving the recurrences gives $u_r\le4\cdot6^{\rho-1-r}$. For the sizes, $\kappa_r=\mu_r/(db)$ satisfies $\kappa_{r-1}\le10\kappa_r+2^{1-r}$
up to the terms per interval, so
\[
  \kappa_r\le\sum_{j=r}^{\rho-1}10^{j-r}2^{1-j}=2^{1-r}\sum_{j=r}^{\rho-1}5^{j-r}\le\tfrac54\,2^{1-r}\,5^{\rho-r},
\]
a bound geometric in $r$; the terms per interval add $O(\sum_j10^{j-r}6^{\rho-j})=O(10^{\rho-r})$ cells at level $r$. Hence
$\sum_{r<\rho}\mu_r\le c\,(5^\rho db+10^\rho)$, $\sum_{r<\rho}\mu_r/h_r=\sum_r\mu_r4^r/L\le c\,(5^\rho/d+40^\rho/L)$, and the backward
stage has at most $c\,(6^\rho+5^\rho/d+40^\rho/L)$ runs.

\paragraph{Rewind and emit.} The rewind moves only tape~$1$, down from $\mathrm{CB}_1$ offset $\mathrm{sh}_1$. It reads the live tainted
intervals of tape~$1$: the comb regions and the scratch windows of levels below $\rho$, and the loaded copy $[db,3db]$; so it has at most
$c\,6^\rho+2\rho+2$ runs. Afterwards only the loaded copy among them is visited again. The emit's step $u$ reads $\mathrm{OUT}_1$
offset $u$, tape-$1$ cell $u$ and tape-$2$ cell $\min(u,\sum_rm_r-1)$, and appends output bit $u$. A tainted step of the emit reads a
tainted cell of the loaded copy, at $u\in[db,3db]$; of $\mathrm{OUT}_1$, which has at most $2u_1+2$ intervals and size at most $4\mu_1$;
or of $Z_r$ for $r<\rho$, which is tainted on $U_r$. Among the outputs $u\in[db,3db]$, those of $\Out(A)$ are the $4b$ bits of the
entries $(1,1)$, $(1,2)$, $(2,1)$ and $(2,2)$. Elsewhere, the outputs of $\Out(A)$ form one run of $2b$ positions in each block of $db$
output positions, so an interval of length $\lambda$ contains at most $2b(\lambda/(db)+2)$ of them. The other tainted intervals of the
emit number at most $c\,6^\rho$ and have total length at most $4\mu_1+\sum_r\mu_r\le c\,(5^\rho db+10^\rho)$. Hence
\[
  N_{\mathrm{up}}\le4b+2b\bigl(c\,5^\rho+c\,10^\rho/(db)+2c\,6^\rho\bigr)\le4b+c'\,6^\rho b\qquad\text{once }(5/3)^\rho\le db .
\]
Altogether, for every $1\le\rho\le r_{\max}$,
\begin{equation}\label{eq:c3-total}
  \abs{C_0}\le\frac{db}{2^\rho}+2,\qquad N_{\mathrm{run}}\le c''\Bigl(6^\rho+\frac{5^\rho}d+\frac{40^\rho}L+\rho\Bigr),\qquad
  N_{\mathrm{up}}\le4b+c'\,6^\rho b .
\end{equation}

\subsection{Cascade 4}\label{app:cascade4}

Cascade~4 uses the closure of order $h=2$, because its final sweep moves tapes~$0$ and~$1$ once every three steps. All its regions are
disjoint, separated by two buffer cells. On tape~$2$: the frames $Z_r=[z_r,z_r+m_r)$ and $\mathrm{BS}_r=[z_r+m_r+2,z_r+3m_r+2)$ for
$r=1,\dots,r_{\max}$, with $z_1=0$ and $z_{r+1}=z_r+3m_r+4$, ending at $E_2=\sum_{r\le r_{\max}}(3m_r+4)<3L$; then the seed region
$\mathrm{SD}=[E_2+2,E_2+2+L)$. On tape~$0$: the loaded copy $\mathrm{LD}_0=[0,L)$; $X_1=[L+2,2L+2)$; for $r=1,\dots,r_{\max}$, a region
$\mathrm{H2R}_r$ of size $m_r$ and the region $X_{r+1}$ of size $m_r=2m_{r+1}$, the input of level $r+1$; then, for $r=r_{\max}$ down to
$1$, regions $\mathrm{CM}_r$ of size $m_r$, $\mathrm{U2R}_r$ of size $2m_r$ and $B_{r-1}$ of size $2m_r$, with $B_{r_{\max}}=X_{r_{\max}+1}$
and $B_0$ of size $L$. On tape~$1$: the loaded copy $\mathrm{LD}_1=[0,L)$; for $r=1,\dots,r_{\max}$, regions $\mathrm{H1}_r$ and
$\mathrm{H2}_r$ of size $m_r$; then, for $r=r_{\max}$ down to $1$, regions $\mathrm{MZ}_r$ of size $m_r$ and $\mathrm{U1}_r$, $\mathrm{U2}_r$
of size $2m_r$. In the table, $Y[x]$ is offset $x$ of region $Y$.

\begin{center}\small
\begin{longtable}{>{\raggedright\arraybackslash}p{0.2\textwidth}>{\raggedright\arraybackslash}p{0.74\textwidth}}
\toprule
phase & head motion \\
\midrule
\endhead
initial positioning, load, rev & tape~$2$ up from $0$ to $\mathrm{SD}[0]$ while the input head stays on bit~$0$; load, step $u<L$:
input head on $u$, tapes $0$ and $1$ on $\mathrm{LD}_0[u]$ and $\mathrm{LD}_1[u]$, tape~$2$ on $\mathrm{SD}[\lceil u/2\rceil]$ while
$u\le\lfloor L/3\rfloor$, then at speed~$1$; tape~$0$ up to $X_1[L-1]$; rev: $X_1[x+\Delta]$ aligned with $\mathrm{SD}[x]$, then tape~$2$
parks on $\mathrm{SD}[0]$; tape~$2$ down to $z_1=0$ \\
fwd $r$.move, transfer & tape~$0$ idle on $X_r[0]$, tape~$1$ up to $\mathrm{H1}_r[0]$; transfer, step $s$: $X_r[\min(s,m_r-1)]$ aligned
with $\mathrm{H1}_r[\lceil s/2\rceil]$ for $s<2\lfloor m_r/3\rfloor$, then at speed~$1$, ending with one step at which no head moves \\
fwd $r$.shift, rev2 & tape~$0$ reads $X_r[m_r],\dots,X_r[2m_r-1]$, cyclically shifted by $\mathrm{sh}_r$, into the fresh $\mathrm{H2}_r$;
then $\mathrm{H2}_r$ is reversed into the fresh $\mathrm{H2R}_r$; tape~$2$ idle on $z_r$ \\
fwd $r$.merge.move, merge, back, move5 & merge.move, $m_r+2$ steps: tape~$0$ down $\mathrm{H2R}_r$ from offset $m_r-1$ to $0$ ($m_r-1$
steps, then it parks), tape~$1$ down from $\mathrm{H2}_r[0]$ to $\mathrm{H1}_r[0]$, so that $\mathrm{H1}_r[x]$ is aligned with
$\mathrm{H2R}_r[x-3]$, tape~$2$ idle on $z_r$; merge: $\mathrm{H2R}_r[u]$ and $\mathrm{H1}_r[u]$ into $Z_r[u]$; back: $Z_r[u]$ into
$X_{r+1}[u]$, downwards; move5 ($r<r_{\max}$): tape~$2$ up from $z_r$ to $z_{r+1}$ across $Z_r$ and the clean $\mathrm{BS}_r$ \\
comb $r$.1a.move, 1a & 1a.move: tape~$0$ up $B_r[0..m_r-1]$; tape~$1$ up across $\mathrm{U2}_{r+1}$ (or $\mathrm{H2}_{r_{\max}}$) and
the fresh $\mathrm{MZ}_r$; tape~$2$ down from $\mathrm{BS}_{r+1}[0]$ across $Z_{r+1}$ and $\mathrm{BS}_r$ to $Z_r[m_r-1]$ (up from
$Z_{r_{\max}}[0]$ if $r=r_{\max}$); 1a: $B_r[u]$ and $Z_r[u]$ into $\mathrm{MZ}_r[u]$, downwards \\
comb $r$.blocks.move, blocks & blocks.move: tape~$0$ up across $B_r$ into the fresh $\mathrm{CM}_r[m_r-1]$, tape~$1$ up $\mathrm{MZ}_r$;
blocks: $2^r$ blocks of height $h_r=m_r/2^r$, from the top; in each, tape~$1$ down $h_r-1$ cells, up $h_r-1$ cells, then tapes $0$
and $1$ down $h_r$ cells together \\
comb $r$.dbl, shift, rev2 & dbl: $\mathrm{CM}_r[z]$ into $\mathrm{U1}_r[2z]$ and $\mathrm{U1}_r[2z+1]$; shift:
$\mathrm{CM}_r[(z+\mathrm{sh}_r)\bmod m_r]$ into $\mathrm{U2}_r[2z]$ and $\mathrm{U2}_r[2z+1]$; rev2: $\mathrm{U2}_r$ reversed into
$\mathrm{U2R}_r$; tape~$2$ idle on $z_r$ \\
comb $r$.merge.move, merge, out & merge.move, $2m_r+2$ steps: tape~$0$ down $\mathrm{U2R}_r$, tape~$1$ down from $\mathrm{U2}_r[0]$ to
$\mathrm{U1}_r[0]$, tape~$2$ up from $Z_r[0]$ to $\mathrm{BS}_r[0]$ ($m_r+2$ steps); merge: $\mathrm{U2R}_r[u]$ and $\mathrm{U1}_r[u]$ into
$\mathrm{BS}_r[u]$; out: $\mathrm{BS}_r[u]$ into $B_{r-1}[u]$, downwards \\
rewind, emit & rewind: tape~$0$ idle on $B_0[0]$, tape~$1$ down from $\mathrm{U1}_1[L-1]$ to cell~$0$, tape~$2$ down from
$\mathrm{BS}_1[0]$ to cell~$0$; emit, steps $3u+i$ for $u<L$ and $i=0,1,2$: tape~$0$ on $B_0[u]$ and tape~$1$ on $\mathrm{LD}_1[u]$, both
moving at $i=2$, tape~$2$ on cell $\min(3u+i,E_2-1)$, output bit $u$ appended at $i=2$ \\
\bottomrule
\end{longtable}
\end{center}

Split each phase into segments at the steps where a head stops early or where no head moves, and each block into its sweeps; in
particular split fwd $r$.transfer where tape~$0$ stops while tape~$1$ still moves, and the emit where tape~$2$ parks on cell
$E_2-1$. Then in every segment each head is idle or monotone, waiting at most two consecutive steps on a cell, and whenever more than one
moving head waits, which happens only in the emit, the moving heads that wait all wait at the same steps.

\begin{lemma}[Segments of Cascade~4]\label{lem:c4-segments}
\notformalized
In a segment $[t_0,t_1)$ of Cascade~4, the number of runs is at most the number of maximal intervals of live tainted cells read by the
moving heads, plus the number of idle heads on a live tainted cell that arrived there at most two steps before $t_0$, plus the number
of runs of reads of $A$; and each run writes at most one interval of cells on each tape whose head moves.
\end{lemma}

\begin{informalproof}
As in \cref{lem:c3-segments}, a tainted step that neither reads a live tainted cell with a moving head nor is one of the first steps
of an idle head on a live tainted cell nor reads $A$ is tainted through a short wait from the step before it. An idle head that parked
three or more steps before $t_0$ has wait age at least $3$ throughout the segment and passes no taint. The moving heads that wait do so
at the same steps, at most two in a row, and then all advance together with the heads that do not wait; so a chain of short waits ends
within two steps. The one longer chain, of at most five tainted steps, occurs when an idle head's short waits at $t_0+1,t_0+2$ hand over
to a moving head's waits at $t_0+3,t_0+4$, at the step where tape~$2$ parks in the emit; it stays adjacent to its run. So each extension
adds no run and only boundedly many written cells.
\end{informalproof}

The cut is placed at a level $\rho$ with $1\le\rho\le r_{\max}$: $C_0$ is the set of tape-$1$ edges that leave tainted steps of
fwd $\rho$.transfer and whose next visit lies after the transfer; as for Cascade~3, this is well defined.

\begin{lemma}[Forward windows of Cascade~4]\label{lem:c4-forward}
\notformalized
For $1\le r\le\rho$ let $W_r$ be the set of offsets of $X_r$ that are live tainted when fwd $r$.transfer starts, and put
$S_r=\{\lceil s/2\rceil:s\in W_r\}$. Then $W_1=\{\lceil u/2\rceil+\Delta:u\in A\}$ and $W_{r+1}=S_r\cup(S_r-3)$ for $r<\rho$. Each $W_r$
is an interval of at least $6$ offsets with
\[
  \abs{W_r}\le\frac{db}{2^{r-1}}+8,\qquad \frac{m_r}3-6\le\min W_r\le\max W_r\le\frac{m_r}3+\frac{3db}{2^r}+1\le2\lfloor m_r/3\rfloor-3 ,
\]
and $S_r\subseteq[7,\mathrm{sh}_r-3]$. Moreover $\abs{C_0}=\abs{S_\rho}\le db/2^\rho+5$, and the forward stage up to the end of
fwd $\rho$.transfer has $O(\rho)$ runs.
\end{lemma}

\begin{informalproof}
Load writes the tainted cells $\mathrm{SD}[\lceil u/2\rceil]$, $u\in A$, since $3db\le\lfloor L/3\rfloor$, together with the loaded
copy $\mathrm{LD}_1[db..3db]$; rev aligns $\mathrm{SD}[x]$ with $X_1[x+\Delta]$, with no waits of the moving heads and tape~$1$ parked for
more than two steps, which gives $W_1$. The bounds pass from $W_r$ to $W_{r+1}$ as in \cref{lem:c3-forward}. If $K$ is an interval of
at least $3$ offsets, $(K-3)\cup K$ is the interval $[\min K-3,\max K]$; since $\abs{S_r}\ge\abs{W_r}/2\ge3$, each $W_{r+1}$ is again one
interval, of at least $6$ offsets, with $\abs{W_{r+1}}\le\abs{W_r}/2+4$, $\min W_{r+1}=\min S_r-3\ge m_{r+1}/3-6$ and
$\max W_{r+1}=\max S_r\le m_{r+1}/3+3db/2^{r+1}+1$. As in \eqref{eq:c3-window}, $\max W_r\le\frac{41}{96}m_r+1\le2\lfloor m_r/3\rfloor-3$,
and $S_r\subseteq[\frac{m_r}6-3,\frac{41}{192}m_r+1]\subseteq[7,\mathrm{sh}_r-3]$, since $m_r\ge64$.

So fwd $r$.transfer stays in its half-speed part on $W_r$ and writes the tainted offsets $S_r$ of $\mathrm{H1}_r$; for $r=\rho$ the
edges leaving them are $C_0$. As $W_r\subseteq[0,m_r-2]$, the second half of $X_r$ is untainted, so shift and rev2 have no tainted step,
and $\mathrm{H2}_r$ and $\mathrm{H2R}_r$ are clean. In fwd $r$.merge.move, tape~$1$ visits $\mathrm{H1}_r[x]$, $x\in S_r$, at the same
step at which tape~$0$ visits $\mathrm{H2R}_r[x-3]$, with no waits of either: tape~$0$ parks on $\mathrm{H2R}_r[0]$ only at
step $m_r-1$ of merge.move, when tape~$1$ is at offset $3$ of $\mathrm{H1}_r$, below $\min S_r\ge7$. Tape~$2$ has been parked on $z_r$ since
before the transfer. So merge taints $Z_r$ exactly at the offsets $S_r\cup(S_r-3)$, and back copies $Z_r$ onto $X_{r+1}$, which gives $W_{r+1}$;
move5 crosses this window and leaves $\mathrm{BS}_r$ untainted. Each forward level has a bounded number of segments, each reading at
most two intervals, which gives the $O(\rho)$ runs.
\end{informalproof}

\begin{lemma}[Clean middle of Cascade~4]\label{lem:c4-middle}
\notformalized
No step of fwd $\rho$ after its transfer, of fwd $r$ for $r>\rho$, or of comb $r$ for $r\ge\rho$ is tainted.
\end{lemma}

\begin{informalproof}
By \cref{lem:c4-forward} the second half of $X_\rho$ and $\mathrm{H1}_\rho[m_\rho-1]$ are untainted, and $C_0$ cuts every tape-$1$ edge
from a tainted step of fwd $\rho$.transfer to the later phases. From the end of that transfer to the end of comb $\rho$.out, tape~$0$
visits only cells at or beyond $X_\rho[m_\rho-1]$, tape~$1$ only cells at or beyond $\mathrm{H1}_\rho$, and tape~$2$ only cells of
$[z_\rho,E_2)$. The live tainted regions of levels at most $\rho$ (on tape~$0$, $\mathrm{LD}_0$, the first halves of $X_{r'}$ for
$r'\le\rho$ and $\mathrm{H2R}_{r'}$ for $r'<\rho$; on tape~$1$, $\mathrm{LD}_1$, $\mathrm{H1}_{r'}$ for $r'<\rho$ and the cut cells of
$\mathrm{H1}_\rho$; on tape~$2$, $\mathrm{SD}$ and $Z_{r'}$ for $r'<\rho$) are not reached by an uncut edge in these phases.
\end{informalproof}

For $1\le r\le\rho-1$ let $U_r\subseteq[0,m_r)$ be the set of offsets at which $B_r$ or $Z_r$ is live tainted when comb $r$.1a starts,
with $u_r$ maximal intervals and $\mu_r=\abs{U_r}$, and let $V_r\subseteq[0,2m_r)$ be the set of live tainted offsets of $\mathrm{BS}_r$
at that moment, with $v_r$ maximal intervals and $\eta_r=\abs{V_r}$. For $r\le\rho-1\le k-2$ the block height is
$h_r=m_r/2^r=d^2b/4^r\ge16b\ge16$.

\begin{lemma}[Backward stage of Cascade~4]\label{lem:c4-backward}
\notformalized
Let $\rho\ge2$. For $1\le r\le\rho-1$, $V_r\subseteq[\frac32m_r+2,2m_r)$; $u_{\rho-1}\le1$, $\mu_{\rho-1}\le db/2^{\rho-1}+8$ and
$V_{\rho-1}=\varnothing$; and for $1<r\le\rho-1$
\begin{gather*}
  u_{r-1}\le10u_r+v_r+19,\qquad v_{r-1}\le7u_r+v_r+13,\\
  \mu_{r-1}\le17\mu_r+\eta_r+\frac{db}{2^{r-1}}+28u_r+30,\qquad \eta_{r-1}\le12\mu_r+\eta_r+20u_r+12 .
\end{gather*}
Comb $r$ has $O(u_r+v_r+u_{r+1}+v_{r+1}+\mu_r/h_r+1)$ runs, with $u_\rho=v_\rho=0$, and leaves $Z_r$ and $\mathrm{BS}_r$ with
$O(u_r+v_r+1)$ tainted intervals and $O(\mu_r+\eta_r+u_r+1)$ tainted cells.
\end{lemma}

\begin{informalproof}
At $r=\rho-1$, comb $\rho$ is clean (\cref{lem:c4-middle}), so comb $(\rho-1)$.1a.move crosses only clean cells and has no tainted step;
$U_{\rho-1}$ is the window $W_\rho$ that merge wrote on $Z_{\rho-1}$, and $V_{\rho-1}=\varnothing$. Now follow comb $r$ for $r\le\rho-1$,
assuming $V_r\subseteq[\frac32m_r+2,2m_r)$.
\begin{enumerate}[label=(\arabic*)]
\item 1a writes a subset of $U_r$ onto the fresh $\mathrm{MZ}_r$. In blocks.move tape~$0$ crosses $B_r$ and tape~$1$ crosses
$\mathrm{MZ}_r$; a short wait of tape~$2$ on $Z_r[0]$ adds at most offset $1$. Tape~$0$ then crosses two clean cells into the fresh
$\mathrm{CM}_r$ while tapes $1$ and~$2$ have wait age at least $3$, so $\mathrm{CM}_r$ enters blocks clean.
\item In each block $[B,B+h_r)$ tape~$1$ turns on $B$ without waiting, and the third sweep ends by stepping onto the next block with
wait age $0$ on both moving heads, so no taint spills across blocks. If $B+h_r-1\in U_r$, the two-step short wait of tape~$0$ at the top
of the block taints at most $\mathrm{MZ}_r[B+h_r-2]$ and $\mathrm{MZ}_r[B+h_r-3]$; otherwise tape~$0$ reaches the third sweep with wait
age $2h_r-2>2$. So the tainted cells of $\mathrm{CM}_r$ form at most $u_r$ intervals with at most $\mu_r+2u_r$ cells, and blocks has
$O(\mu_r/h_r+u_r)$ runs.
\item Dbl writes $\mathrm{U1}_r$ with at most $u_r$ intervals and $2\mu_r+4u_r$ cells; shift and rev2 write $\mathrm{U2}_r$ and
$\mathrm{U2R}_r$ with at most $u_r+3$ intervals and $2\mu_r+4u_r+2$ cells.
\item In merge.move, tapes $0$, $1$ and $2$ cross $\mathrm{U2R}_r$, $\mathrm{U1}_r$ and $Z_r$ together, and tape~$2$ then parks on the
clean $\mathrm{BS}_r[0]$, as $0\notin V_r$. Afterwards $Z_r$ has at most $3u_r+4$ tainted intervals and $5\mu_r+8u_r+3$ tainted cells.
Merge unites $\mathrm{U2R}_r$, $\mathrm{U1}_r$ and $V_r$ on $\mathrm{BS}_r$, and out copies $\mathrm{BS}_r$ onto $B_{r-1}$: at most
$6u_r+v_r+9$ intervals and $10\mu_r+\eta_r+16u_r+8$ cells on each.
\item For $r>1$, in comb $(r-1)$.1a.move, $5m_r+7$ steps, tape~$0$ crosses $B_{r-1}[0..2m_r-1]$, tape~$1$ crosses $\mathrm{U2}_r$ and
enters the fresh $\mathrm{MZ}_{r-1}$, and tape~$2$ crosses two buffer cells, then $Z_r$, and enters the clean $\mathrm{BS}_{r-1}$ at step
$m_r+5$. The buffer cells $z_r+m_r$ and $z_r+m_r+1$ were also crossed by tape~$2$ in comb $r$.merge.move, so they may be tainted; they
meet $B_{r-1}[1..2]$ and add at most one interval and two cells, which the constants $19$ and $30$ of the lemma absorb. From step $2m_r+3$ on, tape~$0$ has wait age at least $4$, and tapes $1$ and $2$ scan only clean cells. So $\mathrm{MZ}_{r-1}$
stays clean, and $\mathrm{BS}_{r-1}$ is tainted only during the steps $[m_r+5,2m_r+2]$, which reach its offsets in
$[3m_r+2,4m_r)=[\frac32m_{r-1}+2,2m_{r-1})$. Counting the intervals and cells read in these steps gives the bounds on $v_{r-1}$ and
$\eta_{r-1}$; counting those read in the steps $[0,2m_r-1]$, together with the window $W_r$ of $Z_{r-1}$ (one interval, at most
$db/2^{r-1}+8$ cells), gives the bounds on $u_{r-1}$ and $\mu_{r-1}$.\qedhere
\end{enumerate}
\end{informalproof}

To solve the recurrences, put $\hat u_r=7u_r+v_r$ and $\hat\mu_r=12\mu_r+\eta_r$. Then $\hat u_{\rho-1}\le7$ and
$\hat u_{r-1}\le77u_r+8v_r+146\le11\hat u_r+146$, so $\hat u_r+14.6\le21.6\cdot11^{\rho-1-r}$ by induction, and
$\hat u_r\le22\cdot11^{\rho-1-r}$. Likewise $\hat\mu_{\rho-1}\le12db/2^{\rho-1}+96$ and
$\hat\mu_{r-1}\le216\mu_r+13\eta_r+12db/2^{r-1}+356u_r+372\le18\hat\mu_r+12db/2^{r-1}+51\hat u_r+372$, whence
\[
  \hat\mu_r\le c\,\bigl(2^{-r}9^{\rho-r}db+18^{\rho-r}\bigr),\qquad \sum_{r<\rho}(\mu_r+\eta_r)\le c\,(9^\rho db+18^\rho),\qquad
  \sum_{r<\rho}\frac{\mu_r}{h_r}\le c\,\Bigl(\frac{9^\rho}d+\frac{18^\rho}{d^2b}\Bigr),
\]
and the backward stage has $O(11^\rho+9^\rho/d+18^\rho/L)$ runs.

\paragraph{Rewind and emit.} In the rewind, tape~$0$ stays on $B_0[0]$ while tapes $1$ and $2$ move down to cell $0$, in two segments
with $O(11^\rho+\rho)$ runs. Tape~$1$ crosses $\mathrm{U1}_1$ first; when it later reaches $\mathrm{LD}_1$, tapes $0$ and $2$ have wait age
far above $2$, so the taint of $\mathrm{LD}_1$ stays in $[db,3db]$. In the emit, the step that appends output $u$ is tainted only if
$B_0[u]$ or $\mathrm{LD}_1[u]$ is live tainted, if one of the tape-$2$ cells $3u,3u+1,3u+2$ below $E_2$ is live tainted, or if
$u=\lceil E_2/3\rceil$ and $B_0[u-1]$ is live tainted (the park of tape~$2$ can extend the tainted interval of $u-1$ by one output).
Within $[db,3db]$ the outputs of $\Out(A)$ are the $4b$ bits of the entries $(1,1)$, $(1,2)$, $(2,1)$ and $(2,2)$. Elsewhere, $B_0$ and the
tape-$2$ frames $Z_r$, $\mathrm{BS}_r$ of levels $r<\rho$ give $O(11^\rho)$ tainted intervals of outputs, of total length
$O(9^\rho db+18^\rho)$, each block of $db$ output positions contains at most $2b$ outputs of $\Out(A)$, and $(18/11)^\rho<2^\rho\le db$.
So
\begin{equation}\label{eq:c4-total}
  \abs{C_0}\le\frac{db}{2^\rho}+5,\qquad N_{\mathrm{run}}\le c''\Bigl(11^\rho+\frac{9^\rho}d+\frac{18^\rho}L+\rho\Bigr),\qquad
  N_{\mathrm{up}}\le4b+c'\,11^\rho b .
\end{equation}

\subsection{Proof of \texorpdfstring{\cref{thm:cascade-cuts}}{the theorem} and \texorpdfstring{\cref{rem:cascade-widths}}{the remark}}

Let $b=b(d)=k$, let $\nu$ be as in \cref{thm:cascade-cuts}, and choose $\rho$ with $\nu\le2^\rho<2\nu$; for $d\ge d_0$ this gives
$1\le\rho\le r_{\max}$ and $(5/3)^\rho\le db$. For Cascade~3, $6^\rho<6\nu^{\log_26}\le6(db)^{0.52}$,
$5^\rho/d\le5(db)^{0.47}/d$ and $40^\rho/L\le40(db)^{1.07}/(d^2b)$, and the last two tend to $0$. For Cascade~4,
$11^\rho<11\nu^{\log_211}\le11(db)^{0.58}$, $9^\rho/d\le9(db)^{0.53}/d$ and $18^\rho/L\le18(db)^{0.7}/(d^2b)$, and again the last two
tend to $0$. Since $b=\log_2d$, \eqref{eq:c3-total} and \eqref{eq:c4-total} give $\abs{C_0}$, $N_{\mathrm{run}}$, $N_{\mathrm{up}}\le c_1\,db/\nu$
for $d\ge d_0$, with an absolute constant $c_1$.

(a) By \cref{lem:closure}(b), $\Gamma$ separates $\Out(\Qstar)=\Out(A)$ from $\Lin(\{1,2\})$, and $\abs\Gamma\le3c_1\,db/\nu$. \Cref{lem:R}
for the oblivious machine ($\eta=0$) gives $(\wval+1)\abs\Gamma\ge\wval\abs\Gamma+\abs{\Gamma\cap\Out(\Qstar)}\ge2db$, so
$\wval+1\ge2\nu/(3c_1)$.

(b) Let $V$ be the set of $v\in A$ whose root $\Out(v)$ is untainted, so $\abs V\ge2db-N_{\mathrm{up}}\ge db$ for $d\ge d_0$. By
\cref{lem:closure}(a), every path from a read of $A$ to a step of $\Out(V)$ uses an edge of $C$, which has at most $\abs{C_0}+3N_{\mathrm{run}}$
cell edges and at most $N_{\mathrm{run}}$ state edges. \Cref{lem:cellcut} gives
$db\le4c_1(db/\nu)\bigl(\log_2(g+1)+\log_2(q+1)\bigr)$, so $\log_2(g+1)+\log_2(q+1)\ge\nu/(4c_1)$.

For \cref{rem:cascade-widths}, let $b\ge1$ be arbitrary and choose $d^{1/5}\le2^\rho<2d^{1/5}$ (Cascade~3) or $d^{1/6}\le2^\rho<2d^{1/6}$
(Cascade~4). The same estimates give $N_{\mathrm{run}}=O(d^{0.52})$, respectively $O(d^{0.58})$, and $N_{\mathrm{up}}\le4b+O(d^{0.52}b)$,
respectively $4b+O(d^{0.58}b)$, uniformly in $b$, while $\abs{C_0}\le d^{4/5}b+2$, respectively $d^{5/6}b+5$. So $\abs V\ge db$ for $d\ge d_0$,
and the proof of \cref{lem:cellcut}, which uses nothing about the entry width, excludes the trace once $\log_2(g+1)+\log_2(q+1)$ is below
$d^{1/5}/c$, respectively $d^{1/6}/c$.


\bibliographystyle{plainurl}
\bibliography{references}

\end{document}
